% !TEX root = ../print.tex
% Draft \S9 --- The subordination bridge
%
% Source: draft/spatial-hemigroup-scale-space.md, \S9 (lines 521-583).
%
% This is the chapter where the causal theory is MATHEMATICS rather than provenance, so it
% is the one place -- besides rem:provenance-causal -- where the causal article is cited
% inside statements and proofs. Its cone is recorded here as def:causal-admissible so that
% the blueprint is self-contained.
%
% REROUTE recorded here: the draft's analytic proof of lem:bridge-exponents cites the
% composition rule for Bernstein functions and continuous negative definite functions
% (Jacob Vol. I, \S3.9). The proof below does not: it goes through Bernstein--Widder
% (ledger A11) and the complete monotonicity of the exponential of a Bernstein function
% (ledger A12), which realises e^{-tau g(omega^2/2)} as a positive mixture of Gaussian
% transforms and so makes it positive definite by def:positive-definite directly. Two small
% cited facts replace one large one, and both are already needed in Chapter 10.
%
% CHANGED (proof of record 2026-09-10): that reroute is now itself superseded for the
% ADMISSIBILITY half of lem:bridge-exponents, which the formalisation proves directly and
% which spends no ledger entry at all. A11 and A12 survive in this chapter only in the
% probabilistic reading of lem:bridge-exponents -- the sentence that asserts the delay T_1
% exists -- and nothing downstream reads that sentence. See the node.
%
% CHANGED (author's decision 2026-09-10): that sentence is no longer part of a statement node.
% lem:bridge-exponents is narrowed to the admissibility clause and a mixture clause CONDITIONAL
% on a delay law, both machine-checked on Lean core; the existence of the delay is
% rem:bridge-subordination, which cites A11(a) and A12 in prose and is marked as not
% machine-checked. So this chapter now spends NO ledger entry inside a statement node, and
% neither A11 nor A12 has a consumer here.

\chapter{The subordination bridge}
\label{ch:bridge}

The causal theory classifies the time-causal families: in its canonical gauge the kernel at
scale $x$ is the law of $x\,T_1$, where $T_1$ is a random delay with Laplace exponent
$F_{\mathrm I}$ of the form recorded below. This article has classified the
reflection-symmetric families. The two cones are related by the oldest construction in the
subject, Bochner's subordination: run a Brownian motion for a random causal delay, and the
displacement reached is symmetric. This chapter makes the construction a map between the two
theories, family to family, and shows that it is not onto. Throughout, $\BM$ is a standard
Brownian motion, so $\mathbb{E}\cos(\omega \BM_u) = e^{-u\omega^2/2}$, and $g_u$ is its law at
time $u$, with $g_0 = \delta_0$; a different diffusion constant is a change of spatial gauge
and is not tracked.

% additive: the causal cone, recorded here so that the statements of this chapter are
% self-contained. It is a definition, not a claim.
% twin: hcs thm:main-characterization + lem:selfdecomposable-exponents (the class those two
% nodes characterise).
% CHANGED (fidelity review 2026-09-10, R4): the closing sentence is requalified. It used to
% assert, unconditionally, that e^{-F_I} is the Laplace transform of the law of a delay
% T_1 >= 0 -- the same existence fact the author's decision 2 of 2026-09-10 moved out of
% lem:bridge-exponents into rem:bridge-subordination, and the one step ledger A11 clause (a)
% and A12 carry. A definition node may not smuggle it back in: complete monotonicity of
% e^{-F_I} is immediate from the displayed form and is what the sentence now says; the delay
% law is cited at the remark, where it is marked as not machine-checked.
% CHANGED (2026-09-11, draft sync, R1): the display now says that the delay profile is
% nonnegative, k_I : (0,infty) -> [0,infty). It was implied by reading k_I(u) du/u as a \Levy{}
% measure, and the Lean carries it as the field CausalAdmissible.k_nonneg, but the class was
% introduced without it. Nothing narrows: every consumer already had the field. Mirrored in the
% draft's Definition 9.0, which also prints the integrability pair for the first time (R11).
% CHANGED (2026-09-12, phase C2, R152): the \lean tag names the SHARED type. This chapter's
% Lean was a deliberate local mirror of Paper I's SelfDecomposableExponent, marked an import
% candidate since 2026-09-08 with TWINS.md's instruction "import, do not restate"; the type now
% lives in ScaleSpaceCore v0.2.0 as ScaleSpace.CausalAdmissible, with this chapter's field
% layout (the two integrability windows of the display below), Paper I's single-finiteness form
% through ofNeTop, and their equivalence ne_top_iff_windows. Paper V reaches it by
% `export ScaleSpace (CausalAdmissible)`, which keeps the unqualified name and dot notation, but
% an alias is not a declaration and a checker that resolves names would miss it -- so the tag is
% the shared name. The statement below is unchanged.
\begin{definition}[the causal admissible cone]\label{def:causal-admissible}
\lean{ScaleSpace.CausalAdmissible, ScaleSpace.CausalAdmissible.exponent}\leanok
A function $F_{\mathrm I} : [0,\infty) \to [0,\infty)$ is \emph{causally admissible} if
\[
  F_{\mathrm I}(\sigma) = b_0\,\sigma
   + \int_0^\infty \bigl(1 - e^{-\sigma u}\bigr)\,k_{\mathrm I}(u)\,\frac{du}{u},
  \qquad b_0 \ge 0, \quad k_{\mathrm I} : (0,\infty) \to [0,\infty) \ \text{nonincreasing},
\]
with $\int_0^1 k_{\mathrm I}(u)\,du < \infty$ and
$\int_1^\infty k_{\mathrm I}(u)\,u^{-1}du < \infty$. These are exactly the exponents the
time-causal article's characterization theorem attaches to its axioms, in its canonical gauge;
$(b_0, k_{\mathrm I})$ is its data, its drift and its delay profile. Such an
$F_{\mathrm I}$ is a Bernstein function vanishing at the origin, so $e^{-F_{\mathrm I}}$ is
completely monotone with value $1$ there; that it is therefore the Laplace transform of the law
of a random delay $T_1 \ge 0$ --- the subordination reading of this chapter --- is cited at
\ref{rem:bridge-subordination} and is not asserted here.
\statusT\quad\emph{(Vocabulary. No claim is made here about the causal axioms; the equivalence
of this class with them is the content of that article, and this chapter uses only the
displayed form. The closing sentence stops at complete monotonicity because the step beyond it
is an existence claim on \ledger{A11}\,(a) and \ledger{A12}, which the author's decision of
2026-09-10 had already moved out of \ref{lem:bridge-exponents}; nothing in the chapter reads
the delay law, every consumer of the bridge quantifying over it instead.)}
\end{definition}

% draft: Lemma 9.1
% NEW as a statement, though the construction is classical. twin: none --- the causal article
% has no bridge chapter.
% CHANGED (author's decision 2026-09-10): the lemma is NARROWED by dropping the delay-law
% clause. The last sentence used to ASSERT that a delay T_1 with Laplace transform e^{-F_I}
% exists, which is the one step ledger A11(a) and A12 carried, and it was typed as
% Skeleton.bridge_delay_law -- an interface with no consumer, since SpatialLine.bridge_exponents
% was proved without reading it and nothing else in the development spends it. The sentence is
% now conditional on such a delay being given, which is exactly what
% SpatialLine.bridge_exponents_mixture proves and what a specification-as-hypothesis loses
% nothing by (the delay law is unique, being determined by its Laplace transform). Dropping a
% hypothesis strengthens the admissibility clause, so nothing that cited this lemma weakens.
% The existence -- the subordination reading -- is now rem:bridge-subordination below, marked as
% not machine-checked and citing A11(a) with A12 in prose. Skeleton.bridge_delay_law is deleted
% and bridge_exponents's unused hypothesis with it; the node is [T] and \leanok.
\begin{lemma}[the bridge on exponents]\label{lem:bridge-exponents}
\uses{def:causal-admissible,def:positive-definite,lem:profile-integrability,lem:selfdecomposable-exponents,prop:fourier-toolbox,thm:main-characterization}
\lean{SpatialLine.bridge_exponents, SpatialLine.bridge_exponents_mixture}\leanok
Let $F_{\mathrm I}$ be causally admissible with data $(b_0, k_{\mathrm I})$ and put
\begin{equation}\label{eq:bridge}
  \Psi(\omega) := F_{\mathrm I}\bigl(\tfrac12\omega^2\bigr), \qquad \omega \in \RR .
\end{equation}
Then $\Psi$ is admissible in the sense of \ref{thm:main-characterization}, that is of the form
\ref{eq:sd-profile}, with Gaussian coefficient $a = \tfrac12 b_0$ and folded profile
\begin{equation}\label{eq:bridge-profile}
  k(x) = 2x\int_0^\infty g_u(x)\,k_{\mathrm I}(u)\,\frac{du}{u}, \qquad x > 0,
\end{equation}
the Gaussian mixture of the delay catalogue. Moreover, if $T_1 \ge 0$ is a random delay
independent of $\BM$ with $\mathbb{E}e^{-\sigma T_1} = e^{-F_{\mathrm I}(\sigma)}$ for every
$\sigma \ge 0$, then the law of $\BM_{T_1}$ --- the mixture of the Gaussian laws against the law
of $T_1$ --- is a symmetric probability law with cosine transform $e^{-\Psi}$.
\statusT\quad\emph{(Both clauses are machine-checked and the node's route is Lean core; no
ledger entry is spent. The delay $T_1$ enters the second clause as a \emph{hypothesis} and not
as an assertion, and quantifying over it loses nothing, a law on $[0,\infty)$ being determined
by its Laplace transform. That a delay with the prescribed transform exists at all is the
subordination reading of the bridge and is \ref{rem:bridge-subordination}, which is cited and
not machine-checked. \textbf{Changed 2026-09-10 (author's decision).} The lemma used to close
by asserting that existence, and carried ledger A11's clause (a) with A12 for it; the
formalisation then proved the admissibility clause without reading the delay law at all, so the
two entries stood behind a statement nothing consumed. Dropping the clause removes a hypothesis
from the admissibility half and is therefore a strengthening: what a proof cites is an upper
bound on what a statement needs, and here the statement \emph{prescribes} the profile, which is
checked field by field, so no representation theorem is needed to produce one. The one field the
draft is right to single out --- that $k$ in \ref{eq:bridge-profile} is nonincreasing, which is
not visible term by term because $\frac{d}{dx}[x\,g_u(x)]$ changes sign at $x = \sqrt u$ ---
becomes visible term by term after the substitution $u = x^2v$; see the proof.
\textbf{Attribution} (2026-09-19, R194, left by batch S2; the pages were read at the held copies
this session). The construction is classical, and its two halves have different sources.
\ref{eq:bridge-profile} is \sscite{sato2001subordination}'s (3.3), p.~8, in the two-sided
convention: his $k^{\sharp}(x) = x\int_0^\infty(2\pi s)^{-1/2}e^{-x^2/2s}\,k(s)\,ds/s$ is half the
folded $k$ printed here, and his (4.2)--(4.3), p.~9, record the same folding at the origin,
$k^{\sharp}(0{+}) = k(0{+})$ and $c_Y = 2c_Z$. That the image of a self-decomposable delay is
self-decomposable is \sscite{halgreen1979self}, \S3, p.~15, for a pure variance mixture --- where
the ratio $\zeta(t^2/2)/\zeta(a^2t^2/2)$ of Laplace transforms is again one, which is the route of
\ref{rem:bridge-route} --- and \sscite{sato2001subordination}, Theorem~1.1, for a general
self-decomposable subordinator, proved there in \S2 by approximating the delay profile by step
functions. What this lemma does that those do not is to \emph{prescribe} the pair $(a,k)$ and
check it field by field, so that no representation theorem is spent and the second clause
quantifies over the delay law instead of producing it.)}
\end{lemma}

\begin{proof}
\leanok
Write $\nu_2(x) = \int_0^\infty g_u(x)\,k_{\mathrm I}(u)\,du/u$ and $k(x) = 2x\,\nu_2(x)$, and
put $a = \tfrac12 b_0$. We verify that $(a,k)$ is a pair of the kind \ref{eq:sd-profile}
requires and that its exponent is $\Psi$; admissibility is then \ref{thm:main-characterization}.

\emph{$\nu_2$ is finite off the origin.} Fix $x \ne 0$. Below $u = 1$ the mixture integrand is
dominated by a multiple of $k_{\mathrm I}$: from $y^2 \le 4e^y$ at $y = x^2/2u$ we get
$e^{-x^2/2u} \le 16u^2/x^4$, and $u \le \sqrt{2\pi u}$ for $u \le 1$, so
$g_u(x)/u \le 16/x^4$ and the integrand is at most $(16/x^4)\,k_{\mathrm I}(u)$, integrable by
$\int_0^1 k_{\mathrm I} < \infty$. Above $u = 1$, $g_u \le 1$ and the integrand is at most
$k_{\mathrm I}(u)u^{-1}$, integrable by the second condition of \ref{def:causal-admissible}. The
two conditions pay in different places, and the first pays only because of the Gaussian: the
mixture carries $du/u$ where that condition carries $du$, and the missing power of $u$ is
supplied by the flatness of $g_u(x)$ as $u \downarrow 0$ at a fixed $x \ne 0$.

\emph{$k$ is nonincreasing.} This is the clause that is not visible term by term. It becomes
visible in the right variable. Substituting $u = x^2v$, which leaves $du/u$ invariant,
\begin{equation}\label{eq:bridge-scaling}
  x\,\nu_2(x) = \int_0^\infty \frac{e^{-1/2v}}{v\sqrt{2\pi v}}\;k_{\mathrm I}(x^2v)\,dv ,
  \qquad x > 0 .
\end{equation}
The weight no longer involves $x$; the whole dependence sits in $k_{\mathrm I}(x^2v)$, which is
nonincreasing in $x > 0$ for each fixed $v > 0$ because $k_{\mathrm I}$ is nonincreasing. Hence
$k = 2x\nu_2(x)$ is nonincreasing on $(0,\infty)$.

\emph{The \Levy{} condition.} Since $\int_\RR (1 \wedge x^2)\,g_u(x)\,dx$ is at most the total mass
$1$ and at most the variance $u$, Tonelli gives
\[
  \int_0^\infty (1 \wedge x^2)\,k(x)\,\frac{dx}{x}
   = \int_0^\infty\Bigl(\int_\RR (1 \wedge x^2)\,g_u(x)\,dx\Bigr) k_{\mathrm I}(u)\,\frac{du}{u}
   \le \int_0^\infty (1 \wedge u)\,k_{\mathrm I}(u)\,\frac{du}{u} < \infty,
\]
the last integral being the single-integral reading of the two conditions of
\ref{def:causal-admissible}. By \ref{lem:profile-integrability} this is the pair of conditions
\ref{eq:sd-profile} asks of $k$. So the spatial \Levy{} condition \emph{is} the causal
integrability condition, with the Gaussian performing the truncation.

\emph{The exponent.} Since $e^{-u\omega^2/2} = \int\cos(\omega x)\,g_u(x)\,dx$ and
$\int g_u = 1$, we have $1 - e^{-u\omega^2/2} = \int_\RR (1 - \cos\omega x)\,g_u(x)\,dx$, and
Tonelli gives
\[
  a\,\omega^2 + \int_0^\infty (1 - \cos\omega x)\,k(x)\,\frac{dx}{x}
   = \tfrac12 b_0\,\omega^2
     + \int_0^\infty \bigl(1 - e^{-u\omega^2/2}\bigr)\,k_{\mathrm I}(u)\,\frac{du}{u}
   = F_{\mathrm I}\bigl(\tfrac12\omega^2\bigr) = \Psi(\omega),
\]
the folding factor $2$ in $k = 2x\nu_2$ being exactly the factor the half-line reading of an
even integral needs.

\emph{The mixture reading.} Let $T_1 \ge 0$ be as in the statement. The mixture of the Gaussian
laws against the law of $T_1$ is a probability measure, each $g_u$ being one, and it is
symmetric, each $g_u$ being symmetric and symmetry passing through the mixture set by set.
Conditioning on $T_1$ and using $\mathbb{E}\cos(\omega \BM_u) = e^{-u\omega^2/2}$,
\[
  \mathbb{E}\cos(\omega \BM_{T_1}) = \mathbb{E}\,e^{-T_1\omega^2/2}
   = e^{-F_{\mathrm I}(\omega^2/2)} = e^{-\Psi(\omega)},
\]
the middle equality being the hypothesis read at $\sigma = \omega^2/2$. So the law of $\BM_{T_1}$
is a symmetric probability law with cosine transform $e^{-\Psi}$. Nothing here reads where the
delay law comes from: only its Laplace transform is used.
\end{proof}

\begin{remark}[the bridge is subordination]\label{rem:bridge-subordination}
A delay of the kind the lemma's second clause hypothesises always exists, and that is what makes
the bridge an instance of Bochner's subordination rather than a formal identity between
exponents. A causally admissible $F_{\mathrm I}$ is a Bernstein function vanishing at the
origin, so $e^{-F_{\mathrm I}}$ is completely monotone --- ledger A12, Schilling et al.
Thm. 3.7(iii) --- with value $1$ at the origin, and a completely monotone function with value $1$
at the origin is the Laplace transform of a probability law on $[0,\infty)$ --- ledger A11
clause (a), Schilling et al. Thm. 1.4. That law is the law of a random delay
$T_1 \ge 0$, unique because a law on the half-line is determined by its Laplace transform, and
the lemma then says that the symmetric family attached to $\Psi$ is the family of laws of
$\BM_{T_1}$: Brownian motion run for the causal delay. The catalogue formula \ref{eq:bridge-profile}
is the same statement on profiles.

This remark is \emph{not} machine-checked, and it is the only part of the bridge that is not:
the two cited facts are Bernstein's theorem and the complete monotonicity of the exponential of
a Bernstein function, and Mathlib carries no Bernstein-function development at the pinned
version. Nothing in the article depends on the remark --- every consumer of the bridge quantifies
over the delay law instead --- so neither entry is admitted as a Lean axiom, and neither name is
on the trust boundary.
\end{remark}

% CHANGED (2026-09-19, article mirror of module B, R174): the remark is reframed as the paper
% reframed it on 2026-09-17 (paper-b/04-bridge.tex, "to be mirrored"): it is a second route to the
% admissibility clause, stated as such, and not the history of an argument the draft once carried.
% The register pass's finding behind the paper-side rewrite is that "the draft" has no referent for
% a reader of the article. The mathematics is unchanged, the two cited facts keep their ledger
% names in prose, and the sentence saying which of the two routes is machine-checked stays here,
% where the blueprint records it (the paper's copy drops it, the author's decision D5).
\begin{remark}[why the image is self-decomposable, in one line]\label{rem:bridge-route}
There is a second route to the admissibility clause, which produces the profile instead of
checking a prescribed one and shows in one line \emph{why} the image is self-decomposable. For
$c \in (0,1)$ the function $g = F_{\mathrm I} - F_{\mathrm I}(c^2\,\cdot)$ is again a Bernstein
function vanishing at the origin, so $e^{-\tau g}$ is completely monotone for every $\tau > 0$
(A12) and hence the Laplace transform of a probability law $\rho_\tau$ on $[0,\infty)$ (A11);
then
$e^{-\tau(\Psi(\omega) - \Psi(c\omega))} = \int_{[0,\infty)} e^{-u\omega^2/2}\rho_\tau(du)$
is a positive mixture of positive definite functions, so
$\Psi - \Psi(c\,\cdot) \in \NDs$ for every $c$, and condition (1) of
\ref{lem:selfdecomposable-exponents} produces \ref{eq:sd-profile}. The route rests on the two
cited facts of the previous remark; the proof above rests on none, which is why it is the proof
of record. What this route does is \emph{produce} a profile, where the lemma \emph{prescribes}
one and the prescription is checked field by field; that check is machine-checked, and this
remark is not.
\end{remark}

% draft: Proposition 9.2
\begin{proposition}[the bridge on families; the corners; the gauges]\label{prop:bridge-families}
\lean{SpatialLine.bridge_families, SpatialLine.bridge_families_delay,
SpatialLine.bridge_families_gamma, SpatialLine.bridge_families_bessel,
SpatialLine.bridge_families_stable}\leanok
% The forward references to the corners' nodes (prop:matern-exponent, prop:student-t,
% prop:two-members) were removed from \uses on 2026-09-07: they made the dependency graph
% cyclic (those nodes use this one), and plasTeX's graph walk does not terminate on a cycle.
% The identification of the images is a computation done here; the corner nodes cite this one.
% CHANGED (proof of record 2026-09-10, wave 5): clause (2) is rewritten to the machine-checked
% route, which runs in the opposite direction to the printed one, and a \uses edge to
% prop:matern-exponent is added because the checked route consumes that node's integral. The
% printed proof reads the causal exponent gamma log(1+sigma) off the hypothesis and then
% evaluates the profile; the Lean statement hypothesises the causal PROFILE k_I(u) = gamma e^-u,
% so the exponent has to be derived, and the cheap derivation is through the profile rather than
% through a causal Frullani integral: eq:bridge says F_I(omega^2/2) IS the exponent of the
% SDProfile carrying eq:bridge-profile, so once the profile is identified the spatial Frullani
% of prop:matern-exponent(1) finishes. The edge to prop:matern-exponent is acyclic now that that
% node's own edge to this one is gone (see 10-corners.tex).
% CHANGED 2026-09-19 (R177, second external review round, referee B finding 4(i) and referee C
% finding 8): CLAUSE (3) IS CONDITIONAL, as SpatialLine.student_subordinated and
% lem:student-subordinated have been since R160 (2026-09-15) and as prop:student-t(3) has been since
% the same day. The clause read "the causal Bessel family, whose kernels are inverse-gamma laws of
% shape a, maps to the Student-t family with 2a degrees of freedom", which a reader takes as
% naming an instance of this node's hypothesis, and the proof's "self-decomposability of the image
% is lem:bridge-exponents" needs exactly what lem:student-subordinated takes as a hypothesis: that
% the inverse-gamma law is the delay law of a causally admissible exponent. That is ledger A18,
% cited at rem:student-ggc and admitted nowhere; the declaration under the tag,
% bridge_families_bessel, proves the MIXTURE IDENTITY and nothing about a causal family. The clause
% now says which half is unconditional and which is not. Safe direction (a hypothesis made
% explicit); no Lean change. The referee also found the same slide in four running sentences and in
% Table 3's caption; the ones inside this repository's blueprint and paper-b sections are repaired
% with it (the paragraph after prop:thorin-subclass's proof in chapter 10, and the chain display
% and caption in paper-b), and the rest are listed for the summary sweep (batch S4).
\uses{lem:bridge-exponents,prop:fourier-uniqueness,thm:main-characterization,prop:matern-exponent}
\notes{bessel-student-corner}
Let $(\Phi^{\mathrm I}_{x,y})$ be a time-causal family satisfying the causal axioms, with
kernels $\mu^{\mathrm I}_{x,y}$ on $[0,\infty)$ written in its canonical gauge and exponent
$F_{\mathrm I}$. Define, for $0 \le s \le t$,
\begin{equation}\label{eq:bridge-families}
  \mu_{s,t} := \int_{[0,\infty)} g_u\;\mu^{\mathrm I}_{s^2,t^2}(du),
  \qquad \Phis{s}{t}f := \mu_{s,t} * f .
\end{equation}
Then $(\Phis{s}{t})$ satisfies (A1)--(A8) and (ND) with $S_\lambda t = \lambda t$; the
coordinate $t$ is its canonical gauge, and its exponent is $F = F_{\mathrm I}(\cdot^2/2)$. In
particular:
\begin{enumerate}
\item the causal pure delay, $\mu^{\mathrm I}_{x,y} = \delta_{y-x}$, maps to the Gaussian family
of \ref{prop:two-members}, $\mu_{s,t} = g_{t^2-s^2}$; the causal drift ray $b_0\sigma$ maps to
the Gaussian ray $\tfrac12 b_0\omega^2$;
\item the causal Gamma family, $F_{\mathrm I}(\sigma) = \gamma\log(1+\sigma)$, maps to the
\Matern{} family at range $1/\sqrt2$: $F(\omega) = \gamma\log(1 + \omega^2/2)$, folded profile
$k(x) = 2\gamma e^{-\sqrt2\,x}$;
\item if the inverse-gamma laws of shape $a$ are the delay kernels of such a family --- the causal
Bessel family, whose existence is cited at \ref{rem:student-ggc} and asserted nowhere here --- then
it maps to the Student-t family with $2a$ degrees of freedom. The mixture identity behind that
clause is unconditional: mixing the Gaussian kernels against the inverse-gamma law of shape $a$
gives the Student-t law with $2a$ degrees of freedom, which is \ref{prop:student-t}(1);
\item the causal stable family, $F_{\mathrm I}(\sigma) = \sigma^\alpha$ with
$0 < \alpha \le 1$, maps to the symmetric stable family
$F(\omega) = 2^{-\alpha}|\omega|^{2\alpha}$ of index $2\alpha \in (0,2]$; the boundary
$\alpha = 1$ is item (1);
\item the spatial canonical gauge is the square root of the causal one: a causal scale $x$
becomes the spatial scale $t = \sqrt x$, which is the parabolic gauge of the causal article.
\end{enumerate}
\statusT\quad\emph{(Every clause is a computation on exponents once \ref{lem:bridge-exponents}
has supplied admissibility. Item (3)'s density identification is
\ref{prop:student-t}(1). The chain of equalities in item (2) is checked against the profile of
\ref{prop:matern-exponent} at range $1/\sqrt2$ and agrees; see the proof.
\emph{Item (3) is conditional} since 2026-09-19 (R177): what the declarations prove is the
mixture identity, \texttt{SpatialLine.bridge\_families\_bessel}, and the passage from it to a
spatial \emph{family} needs the inverse-gamma law to be the delay law of a causally admissible
exponent --- the hypothesis \ref{lem:student-subordinated} quantifies over and
\ref{rem:student-ggc} cites, ledger \ledger{A18}, admitted nowhere. The other four items name
causal families this article's own computations produce or whose admissibility
\ref{lem:bridge-exponents} supplies from an explicit profile.)}
\end{proposition}

\begin{proof}
\leanok
\emph{The family.} $\hat\mu_{s,t}(\omega) = \int e^{-u\omega^2/2}\,\mu^{\mathrm I}_{s^2,t^2}(du)
= \widehat{\mu^{\mathrm I}_{s^2,t^2}}(\omega^2/2)
= \exp[-(F_{\mathrm I}(t^2\omega^2/2) - F_{\mathrm I}(s^2\omega^2/2))]$, using the causal
similarity form in the canonical gauge. So $g_{s,t}(\omega) = F(t\omega) - F(s\omega)$ with
$F = F_{\mathrm I}(\cdot^2/2)$, admissible by \ref{lem:bridge-exponents}; the family is
therefore the one \ref{thm:main-characterization} attaches to $F$ in the canonical gauge, and
that construction verified (A1)--(A8) and (ND), with $S_\lambda t = \lambda t$ because
$F(\lambda t\omega) - F(\lambda s\omega) = g_{s,t}(\lambda\omega)$. Directly:
$\Dil{\lambda}g_u = g_{\lambda^2u}$, so a spatial dilation by $\lambda$ is a causal dilation by
$\lambda^2$, and the causal covariance
$\Dil{\lambda^2}\mu^{\mathrm I}_{x,y} = \mu^{\mathrm I}_{\lambda^2x,\lambda^2y}$ becomes
$\Dil{\lambda}\mu_{s,t} = \mu_{\lambda s,\lambda t}$; the cascade is inherited from the
independent increments of both $\BM$ and the delay.

Two identifications in that paragraph need justification, and they are where the work is
once \ref{lem:bridge-exponents} is in hand. First, the kernels \ref{thm:main-characterization}
attaches to $F$ and the mixtures \ref{eq:bridge-families} are \emph{a priori} two different
families; they agree because both are symmetric probability measures and their cosine
transforms agree, so \ref{prop:fourier-uniqueness} identifies them. Second, the scaling action
is $S_\lambda t = \lambda t$ on the nose, not merely up to the gauge: the construction of
\ref{thm:main-characterization} produces the gauge conjugate $\chi^{-1}(\lambda\chi(t))$, which
in the canonical gauge $\chi = \mathrm{id}$ is $\lambda t$ for every $t \ge 0$, and every clause
of (A8) reads $S_\lambda$ on $[0,\infty)$ only.

(1) $\mu^{\mathrm I}_{x,y} = \delta_{y-x}$ gives $\mu_{s,t} = g_{t^2-s^2}$, and
$F_{\mathrm I}(\sigma) = b_0\sigma$ gives $F(\omega) = \tfrac12 b_0\omega^2$.

(2) Take the causal Gamma family by its \emph{profile}, $k_{\mathrm I}(u) = \gamma e^{-u}$ with
no drift; the exponent then comes out of the computation rather than into it. By
\ref{eq:bridge-profile} the image has folded profile
\[
  k(x) = 2\gamma x\int_0^\infty g_u(x)\,e^{-u}\,\frac{du}{u}
       = \frac{2\gamma x}{\sqrt{2\pi}}\int_0^\infty u^{-3/2}e^{-x^2/2u}e^{-u}\,du
       = 2\gamma e^{-\sqrt2\,x},
\]
by the identity $\int_0^\infty u^{-3/2}e^{-A/u - Bu}\,du = \sqrt{\pi/A}\,e^{-2\sqrt{AB}}$ at
$A = x^2/2$, $B = 1$ --- equivalently, the Laplace transform
$\int_0^\infty \frac{x}{\sqrt{2\pi u^3}}e^{-x^2/2u}e^{-\sigma u}du = e^{-x\sqrt{2\sigma}}$ of the
first-passage law of Brownian motion at level $x$, at $\sigma = 1$. That is the folded profile
$2\gamma e^{-x/t}$ of \ref{prop:matern-exponent} at $t = 1/\sqrt2$. The exponent now follows
without a second integral: \ref{eq:bridge} says that $F_{\mathrm I}(\omega^2/2)$ \emph{is} the
exponent of the admissible datum whose folded profile is \ref{eq:bridge-profile}, and that
profile has just been identified, so \ref{prop:matern-exponent}(1) evaluates it and
$F(\omega) = \gamma\log(1 + \omega^2/2) = \gamma\log(1 + t^2\omega^2)$ at $t = 1/\sqrt2$. In
particular the causal exponent $\gamma\log(1+\sigma)$ is a consequence of the causal profile and
need not be assumed.

(3) is the classical normal variance-mixture representation, computed in
\ref{prop:student-t}(1); that identity needs no hypothesis on the causal side. Under the clause's
hypothesis --- that the inverse-gamma law of shape $a$ is the delay law of a causally admissible
$F_{\mathrm I}$ --- \ref{lem:bridge-exponents} makes the image admissible and
self-decomposable, and \ref{lem:student-subordinated} is the same statement read as membership in
$\mathcal{S}$.

(4) $(\omega^2/2)^\alpha = 2^{-\alpha}|\omega|^{2\alpha}$, and $\alpha \le 1$ maps onto index
at most $2$.

(5) The causal accumulated exponent in the canonical gauge is $F_{\mathrm I}(x\sigma)$, and the
image has $\hat\mu_{0,t}(\omega) = e^{-F_{\mathrm I}(t^2\omega^2/2)} = e^{-F(t\omega)}$, so the
spatial gauge $t$ satisfies $t^2 = x$.
\end{proof}

The last item deserves a sentence. The causal article names its parabolic gauge because in it
the memory-line evolution of its Bessel family is a Bessel operator, a heat-type generator; the
reason is now visible from the other side, that parameter being the spatial scale of the
subordinated family, with the factor $\tfrac12$ from Brownian scaling. The two articles'
natural coordinates agree through $x = t^2$, causal scale being a variance and spatial scale a
standard deviation --- which is also why the scale-space literature's diffusion variable is the
square of this article's $t$ (\ref{rem:gauge-freedom}).

% draft: Proposition 9.3
% CHANGED (review 2026-09-09, R16): clause (3)'s closing sentence now excludes F = 0. The zero
% exponent is stationary everywhere and IS subordinated -- F_I = 0 is causally admissible -- so
% the sentence was false as written. Clause (2) already carries the same exclusion in the words
% "other than 0", and clause (3) needs it for the same reason. Found by writing the Lean
% statement (Skeleton.bridge_strictness_stationary), whose second conjunct now carries it.
%
% CHANGED (proof 2026-09-10): the node was [A] on ledger A11 and is now [T]. Formalising
% clause (2) showed A11 is not its obligation: def:causal-admissible hands the proof the
% representing measure, so what is left is the positivity of an integral. A11's clause (b) in
% AXIOMS.md is retired with no consumer; the entry stays, spent by prop:thorin-subclass alone.
%
% CHANGED (proof of record 2026-09-10): clauses (2) and (4) are rewritten to the machine-checked
% route. (2) now derives F_I' by differentiation under the integral sign, with the dominating
% function named, and concludes by positivity of the integral instead of by Bernstein; the
% Bernstein reading survives as a closing paragraph, marked as what is not needed. (4) now names
% the work the phrase "image of the causal cone" hid, which is that the causal cone is closed
% under dilation -- the one closure property that is not termwise.
%
% CHANGED (fidelity review 2026-09-10, R76): TWO NARROWINGS AGAINST THE DRAFT ARE RECORDED, both
% of them made silently when the node was written and found by the F8 sweep.
%  (a) The draft's clause (1) closes with a second sentence: "when F is C^1 on (0,infty), a
%      necessary condition is that u |-> F'(sqrt(2u))/sqrt(2u) be completely monotone." The node
%      states the criterion alone. The sentence is TRUE -- F_I' is the Laplace transform of
%      b_0 delta_0 + k_I(u) du, and the chain rule turns it into the displayed map -- but it is
%      strictly weaker than clause (1), which is an iff, and nothing in the chapter or in the
%      Lean uses it: no declaration under the tag mentions complete monotonicity. It is dropped
%      as a redundant weakening, not as a doubtful one.
%  (b) The draft's clause (4) lists "and their Thorin-type superpositions" among the members of
%      S. The node's clause (4) does not. That membership is the image half of
%      prop:thorin-subclass(4) and belongs there; it is also, like the four named families, not
%      machine-checked -- see row R44.
%
% CHANGED (fidelity review 2026-09-10, R44): THE NODE GOES FROM \leanok TO \notready, and the
% statement is untouched. Three passes found the same thing independently -- the supply-chain
% card pass, the F7 import sweep and the F8 draft diff -- so it is recorded once, here. Clause
% (4) has two halves. The CLOSURE half (convex subcone, closed under dilation) and the
% EXCLUSION half (no Cin ray) are bridge_strictness_cone, proved, Lean core. The MEMBERSHIP
% half -- that S contains the Gaussian ray and the Matern, Student-t and symmetric stable
% families -- has no declaration under this tag and none anywhere: IsSubordinated is an
% existential over CausalAdmissible, and the development constructs NO causally admissible
% datum at all, only the add/smul/dilate combinators of CausalCone.lean. prop:bridge-families,
% which the proof cites for the memberships, states the image identities for a GIVEN causal
% datum and exhibits none either.
% Why it is not a cheap fix, which is the part worth writing down: the Gaussian membership is
% one 20-line datum (b_0 = 2a, k = 0) and the causal Gamma one is a middling exercise, but the
% Student-t membership needs the \Levy{} data of the inverse-gamma delay law, which is exactly
% what ledger A18 grounds -- and A18 is deliberately unadmitted, for want of a consumer
% (AXIOMS.md A18: "what the entry grounds is the CAUSAL admissibility of the Bessel family").
% So the membership half cannot be closed without either admitting A18 or splitting the clause.
% Why \notready rather than a disclosure in place: this article's own rule, applied twice
% during the campaign (prop:thorin-subclass clause (3) -> lem:thorin-ggc; prop:polya-frequency's
% fourth exclusion -> lem:polya-student-exclusion), is that a \leanok node claims exactly what
% its declarations prove and an unformalised clause is SPLIT OUT, not disclosed in place. The
% split is a statement change at a Tier-1 node and belongs to the author, so the review takes
% the honest interim: \notready, with the annotation saying precisely which half is checked.
% That interim ended on 2026-09-11; the next marker is what replaced it.
%
% CHANGED (2026-09-11, R44 split): THE MEMBERSHIP SENTENCE IS SPLIT OUT, and the node returns to
% \leanok, statement and proof block both. The author's decision applies the rule the block above
% names: a \leanok node claims exactly what its declarations prove, and an unformalised clause is
% split out rather than disclosed in place (prop:thorin-subclass(3) -> lem:thorin-ggc;
% prop:polya-frequency's fourth exclusion -> lem:polya-student-exclusion). What moved: clause
% (4)'s sentence "containing the Gaussian ray and the Matern, Student-t and symmetric stable
% families" is now lem:subordinated-members below -- [A] on ledger A18, notready, typed as
% Skeleton.subordinated_members -- and the citation of prop:bridge-families that carried it moved
% out of the proof of clause (4) with it. What stays is what the five declarations prove: the
% convex subcone, the dilations, and the exclusion of the Cin rays.
%
% ONE FURTHER NARROWING, found while making the split and NOT part of R44. Clause (4) closed with
% "So the symmetric theory is strictly larger than the causal theory in Brownian disguise, and the
% extra members are exactly those with catalogue mass on hard steps." The first half stays: it is
% the Cin exclusion read against lem:cin-rays, which puts those rays in the admissible cone. The
% second half is an IFF and only one direction of it is proved -- clause (3) gives that an F with
% a lattice-step profile is not subordinated, never the converse -- and the iff is FALSE as it
% stands: an admissible F whose catalogue is carried by a compact subset of (0,infty) and is not
% lattice-carried has no causal ancestor either, every subordinated profile being a Gaussian
% mixture of the delay catalogue (eq:bridge-profile) and hence strictly positive on all of
% (0,infty). So the sentence is not a clause of this node. It is moved into the discussion
% paragraph after the proof and restated in the direction that is proved, and the identification
% that paragraph already carried ("are the members of the spatial theory with no causal
% ancestor") is narrowed with it. The chapter's own CHECK note at rem:bridge-crossings, which
% records that the exact extent of S is open, is the same observation from the other side.
% CHANGED (2026-09-11, R121): the annotation's last sentence no longer says the membership node
% is notready; three of its four memberships are machine-checked there and the Student-t one is
% lem:student-subordinated. Annotation only; the statement and proof are unchanged.
% CHANGED (2026-09-19, article mirror of module B, R174): clause (4)'s closing sentence -- "So the
% symmetric theory is strictly larger than the causal theory in Brownian disguise" -- is MOVED OUT
% of the statement, to the discussion paragraph after lem:student-subordinated's proof, and
% restated without the image ("larger than the image of the causal theory under the bridge"). It is
% the register pass's item 13: commentary inside a statement. It is a reading of the exclusion
% clause and no declaration of the tag proves a sentence of that shape, so the statement claims
% less; the reading survives where the chapter discusses strictness. No Lean and no status change.
\begin{proposition}[strictness: the bridge is not onto]\label{prop:bridge-strictness}
\uses{lem:bridge-exponents,lem:selfdecomposable-exponents,prop:choquet-cone,lem:cin-rays,lem:lattice-zero,def:causal-admissible,prop:bridge-families}
\lean{SpatialLine.bridge_strictness_criterion, SpatialLine.bridge_strictness_increasing,
SpatialLine.bridge_strictness_cin, SpatialLine.bridge_strictness_stationary,
SpatialLine.bridge_strictness_cone}\leanok
\notes{dickman-cin-rays}
Call an admissible $F$ \emph{subordinated} if $F(\omega) = F_{\mathrm I}(\omega^2/2)$ for some
causally admissible $F_{\mathrm I}$, and write $\mathcal{S}$ for the set of subordinated
exponents.
\begin{enumerate}
\item $F$ is subordinated if and only if $\sigma \mapsto F(\sqrt{2\sigma})$ is causally
admissible.
\item Every subordinated $F$ other than $0$ is strictly increasing on $(0,\infty)$, with
$F' > 0$ there.
\item No $\Cin$ ray is subordinated, since $\Cin'(z) = (1 - \cos z)/z$ vanishes at every
$z \in 2\pi\ZZ \setminus \{0\}$. More generally, an admissible $F$ has a stationary point
$\omega_0 \ne 0$ if and only if $a = 0$ and the Choquet measure $\varpi$ of
\ref{prop:choquet-cone} is carried by the lattice $\tfrac{2\pi}{|\omega_0|}\ZZ$, that is, if
and only if $a = 0$ and the profile $k$ is a step function with all its jumps on that lattice; and no such
$F$ with $F \not\equiv 0$ is subordinated.
\item $\mathcal{S}$ is a convex subcone of the admissible cone, closed under the dilations
$F \mapsto F(\lambda\,\cdot)$; it contains none of the extreme rays $\Cin(\tau\,\cdot)$.
\end{enumerate}
\statusT\quad\emph{(The node carried A11 until 2026-09-10, for one step of clause (2): that a
completely monotone function vanishing at an interior point vanishes identically. Formalising
the clause showed that step is not its obligation. \ref{def:causal-admissible} \emph{hands} the
proof the representing measure $k_{\mathrm I}(u)\,du$; what has to be shown is that
$b_0 + \int_0^\infty e^{-\sigma u}k_{\mathrm I}(u)\,du$ is strictly positive unless
$(b_0, k_{\mathrm I})$ vanishes, which is an integral of a positive integrand and needs no
representation theorem. Bernstein's theorem produces a measure; here there is one already. All
five declarations are machine-checked to Lean core, and the proof below is the checked route.
Clause (4)'s \emph{membership} sentence --- that $\mathcal{S}$ contains the Gaussian ray and the
\Matern{}, Student-t and symmetric stable families --- was split out to
\ref{lem:subordinated-members} on 2026-09-11 (R44, the author's decision); its Gaussian, \Matern{}
and stable memberships are machine-checked there, and its Student-t membership is
\ref{lem:student-subordinated}, \texttt{notready} on ledger A18 (R121). What the four clauses
claim here is what the five declarations prove.)}
\end{proposition}

% CHANGED (2026-09-11, R44 split): the proof block is \leanok again. The membership sentence and
% the citation of prop:bridge-families that carried it moved to lem:subordinated-members with the
% clause; what is below is the checked route for all four clauses.
\begin{proof}\leanok
(1) is \ref{def:causal-admissible} read on $F_{\mathrm I}(\sigma) := F(\sqrt{2\sigma})$.

(2) Let $F = F_{\mathrm I}(\cdot^2/2)$ with $F_{\mathrm I}$ causally admissible, and suppose
$F \not\equiv 0$. Fix $\sigma > 0$. On $(1,\infty)$ the elementary bound
$u\,e^{-\sigma u/2} \le 2/\sigma$ turns $e^{-\sigma u/2}k_{\mathrm I}(u)$ into
$(2/\sigma)\,k_{\mathrm I}(u)u^{-1}$, integrable by the second condition of
\ref{def:causal-admissible}, and on $(0,1)$ it is at most $k_{\mathrm I}$, integrable by the
first; so $e^{-\sigma u/2}k_{\mathrm I}(u)$ dominates the $\sigma$-derivative
$e^{-\sigma u}k_{\mathrm I}(u)$ of the integrand uniformly for $\sigma$ in a neighbourhood of
the given one, and differentiation under the integral sign gives
\[
  F_{\mathrm I}'(\sigma) = b_0 + \int_0^\infty e^{-\sigma u}\,k_{\mathrm I}(u)\,du,
  \qquad \sigma > 0 .
\]
Both terms are nonnegative. If the sum vanishes at one $\sigma > 0$ then $b_0 = 0$ and, the
exponential being strictly positive, $k_{\mathrm I} = 0$ almost everywhere on $(0,\infty)$;
but then the representation makes $F_{\mathrm I}$, and so $F$, identically zero. Hence
$F_{\mathrm I}' > 0$ on $(0,\infty)$, and the chain rule gives
$F'(\omega) = \omega\,F_{\mathrm I}'(\omega^2/2) > 0$ for $\omega > 0$. A function continuous
on $(0,\infty)$ with positive derivative there is strictly increasing there, and $F$ is
continuous, its representation \ref{eq:sd-profile} being so by dominated convergence.

The derivative formula is the only place the causal representation is differentiated, and one
classical step does \emph{not} enter it. Read as a function of $\sigma$, $F_{\mathrm I}'$ is
the Laplace transform of the positive measure $b_0\delta_0 + k_{\mathrm I}(u)\,du$, so it is
completely monotone and one may argue, as the causal article does, that a completely monotone
function vanishing at an interior point vanishes identically. That is Bernstein's theorem, and
it is not needed: it \emph{produces} a representing measure, and \ref{def:causal-admissible}
has already given one. What is left is the positivity of an integral, which is the display
above.

(3) $\Cin'(2\pi n) = 0$ for every nonzero integer $n$, which contradicts (2); so no $\Cin$ ray
is subordinated. For the general statement, \ref{eq:symbol} gives
$\omega F'(\omega) = B(\omega) = 2a\omega^2 + \int (1 - \cos\omega u)\,\varpi(du)$, a sum of
nonnegative terms, which vanishes at $\omega_0 \ne 0$ if and only if $a = 0$ and
$\cos(\omega_0 u) = 1$ for $\varpi$-almost every $u$; the latter says, by the argument of
\ref{lem:lattice-zero} applied to $\varpi$, that $\varpi$ is carried by
$\tfrac{2\pi}{|\omega_0|}\ZZ$, which is the profile $k = \varpi((\cdot,\infty))$ being a step
function with jumps on that lattice. Such an $F$ has a stationary point off the origin, so it
is not subordinated, by (2).

(4) $\mathcal{S}$ is the image of the causal cone under $F_{\mathrm I} \mapsto
F_{\mathrm I}(\cdot^2/2)$, so the three closure properties are closure properties of the causal
cone, and the map carries them across because it is linear and because
$F_{\mathrm I}(\lambda^2\,\cdot)$ maps to $F(\lambda\,\cdot)$ --- a spatial dilation by
$\lambda$ is a causal dilation by $\lambda^2$. On the causal side, sums and nonnegative
multiples act on the data $(b_0, k_{\mathrm I})$ termwise and the two conditions of
\ref{def:causal-admissible} are preserved term by term. The dilation $F_{\mathrm I} \mapsto
F_{\mathrm I}(c\,\cdot)$, $c > 0$, has data $(c\,b_0, k_{\mathrm I}(\cdot/c))$, and for it the
two conditions are better read as the single one
\[
  \int_0^\infty (1 \wedge u)\,k_{\mathrm I}(u)\,\frac{du}{u} < \infty,
\]
which they are equivalent to, $(1 \wedge u)/u$ being $1$ below $1$ and $u^{-1}$ above it. The
change of variables $u = cx$ leaves $du/u$ invariant, so it transforms that integral into
$\int_0^\infty (1 \wedge cx)\,k_{\mathrm I}(x)\,dx/x$, and
$1 \wedge cx \le \max(1,c)\,(1 \wedge x)$ bounds it by $\max(1,c)$ times the undilated one. Splitting the two windows first would
instead require the change of variables on two moving intervals. The exclusion of the $\Cin$
rays is (3), whose Choquet coordinates are \ref{prop:choquet-cone}.
\end{proof}

% draft: Proposition 9.3(4), the membership sentence
% NEW as a node (2026-09-11, R44, the author's decision). The content is not new: the sentence
% stood inside prop:bridge-strictness(4) from the first writing of this chapter. What is new is
% that it is stated apart, because it has no declaration and cannot cheaply get one.
% IsSubordinated is an existential over CausalAdmissible, and the development exhibits no causally
% admissible datum at all -- CausalCone.lean has the add/smul/dilate combinators and no generator
% to feed them; prop:bridge-families, which the proof below cites, states the image identities for
% a GIVEN causal datum and exhibits none either. Per family the cost is one datum: Gaussian
% b_0 = 2a with k_I = 0, about twenty lines; the causal Gamma datum for the Matern and the causal
% stable datum for the stable family, middling exercises; and for the Student-t the \Levy{} data of
% the inverse-gamma delay law, which is exactly what ledger A18 grounds and which this article
% leaves unadmitted for want of a consumer. So the node is [A] on A18 and notready, with
% Skeleton.subordinated_members as its typed target. The Lean statement quantifies over the
% Student-t exponent through its transform -- there is no closed-form Student-t exponent in the
% development, and -log of prop:student-t(2)'s closed form is junk at omega = 0 -- which loses
% nothing, exp being injective.
% CARRIED TO THE DRAFT SYNC, NOT STATED HERE (R44's F8 finding): draft Proposition 9.3(4) lists
% "and their Thorin-type superpositions" among the members of S, and the blueprint dropped that
% without a record. It is the image half of prop:thorin-subclass(4) and belongs there; it is not
% a clause of this node, and no declaration of this node would cover it.
% CHANGED (2026-09-11, the author's decision on the post-sync recommendation, R115): the phrase is
% DROPPED, not moved into prop:thorin-subclass(4). That clause already states the bridge on the
% Thorin subclasses, and a membership sentence beside it would repeat the shape R44 and R79
% removed from \leanok nodes; "Thorin-type superpositions" of the four named families is also not
% well defined (the stable family is itself a Thorin superposition). No sentence citing clause (4)
% for "every symmetric Thorin member is subordinated" is added: SpatialLine.thorin_bridge proves
% the forward map only -- the image of a causal Thorin member is a symmetric Thorin member with
% U twice the image of U_I -- and not the surjectivity that clause (4)'s "onto" asserts. The
% draft's bracketed note at Lemma 9.3' is removed.
% CHANGED (2026-09-11, R121, split under the safe-direction delegation): THE THREE ELEMENTARY
% MEMBERSHIPS ARE MACHINE-CHECKED AND THE STUDENT-T MEMBERSHIP IS SPLIT OUT. The statement is
% narrowed to the Gaussian ray and the Matern and symmetric stable families, which is what
% SpatialLine.subordinated_members proves on Lean core; the node is [T] and \leanok and A18 leaves
% it (the author's decision of 2026-09-11: split, not admit). The Student-t membership is
% lem:student-subordinated below, [A] on A18 and notready, with the paragraph of the old proof
% that concerned it; the \uses edge to prop:student-t moved with it. The edge to prop:two-members
% is dropped: no declaration of this node reads it, the families' exponents being the vocabulary
% of prop:matern-exponent and prop:stable-family. The proof is rewritten to the checked route.
% Findings: the development now exhibits causally admissible data
% (SpatialLine/CausalData.lean: the drift, the Gamma profile, the stable profile); and the causal
% stable exponent sigma^beta needs no limit at the origin -- the causal article matches
% derivatives on (0,infty) and limits at 0+, while here the homogeneity of the profile turns the
% causal dilation into a multiple and one derivative at sigma = 1 fixes the constant.
% twin: hcs prop:gamma-family, prop:stable-family -- the causal data
% (Hemigroup.SelfDecomposableExponent.gammaExponent, ...stableExponent), ported onto the local
% mirror CausalAdmissible with the same profiles and guards.
\begin{lemma}[the subordinated slice's members]\label{lem:subordinated-members}
\uses{prop:bridge-strictness,prop:bridge-families,def:causal-admissible,prop:matern-exponent,prop:stable-family}
\lean{SpatialLine.subordinated_members}\leanok
\notes{dickman-cin-rays}
$\mathcal{S}$ contains the Gaussian ray and the \Matern{} and symmetric stable families.
\statusT\quad\emph{(All three memberships are machine-checked to Lean core, as
\texttt{SpatialLine.subordinated\_members}, at the parameter ranges the corner nodes use:
$a \ge 0$ for the Gaussian ray $a\omega^2$, $\gamma, \theta > 0$ for the \Matern{} family, and
$0 < \alpha < 2$ for the symmetric stable family $|\omega|^\alpha$, whose endpoint $\alpha = 2$
is the Gaussian ray. The proof below is the checked route; the three causally admissible data it
exhibits are the first instances of \ref{def:causal-admissible} in the development. The
Student-t family's membership rests on ledger A18 and is stated apart as
\ref{lem:student-subordinated}.
\emph{Provenance.} The sentence stood inside \ref{prop:bridge-strictness}(4) until 2026-09-11,
with no declaration under that node's tag; the fidelity review recorded that as row R44, and the
author's decision split it out on the article's own rule, that a node marked as formalised claims
exactly what its declarations prove. The same day the three elementary memberships were
formalised and the Student-t membership was split out in turn (R121).)}
\end{lemma}

\begin{proof}\leanok
Membership in $\mathcal{S}$ asks for a causally admissible datum $(b_0, k_{\mathrm I})$ whose
exponent, read at $\omega^2/2$, is the given exponent. One datum is exhibited per family, and the
named families at their general parameters are then reached by the closure clauses of
\ref{prop:bridge-strictness}(4).

\emph{The Gaussian ray.} For $b_0 \ge 0$ the datum $(b_0, 0)$ satisfies
\ref{def:causal-admissible} with nothing to check, and its exponent is
$F_{\mathrm I}(\sigma) = b_0\sigma$ for $\sigma \ge 0$. With $b_0 = 2a$ it gives
$F_{\mathrm I}(\omega^2/2) = a\omega^2$.

\emph{The \Matern{} family.} Let $b_0 = 0$ and $k_{\mathrm I}(u) = \gamma e^{-u}$ for $u > 0$. The
profile is nonincreasing, bounded by $\gamma$ on $(0,1)$, and $k_{\mathrm I}(u)/u \le \gamma
e^{-u}$ on $(1,\infty)$, so both conditions of \ref{def:causal-admissible} hold. By
\ref{prop:bridge-families}(2), $F_{\mathrm I}(\omega^2/2) = \gamma\log(1 + \omega^2/2)$, the
\Matern{} exponent at range $1/\sqrt2$. For $\theta > 0$,
\[
  \gamma\log(1 + \theta^2\omega^2) = \gamma\log\bigl(1 + \tfrac12(\sqrt2\,\theta\,\omega)^2\bigr),
\]
so the \Matern{} exponent at range $\theta$ is the dilation by $\lambda = \sqrt2\,\theta$ of the one
at range $1/\sqrt2$, and it lies in $\mathcal{S}$ because $\mathcal{S}$ is closed under
dilations.

\emph{The symmetric stable family.} Let $0 < \beta < 1$ and $c_\beta = \beta/\Gamma(1-\beta)$,
and take $b_0 = 0$ and $k_{\mathrm I}(u) = c_\beta u^{-\beta}$ for $u > 0$. The profile is
nonincreasing, and $\int_0^1 u^{-\beta}du = (1-\beta)^{-1}$ and
$\int_1^\infty u^{-\beta-1}du = \beta^{-1}$ are finite, the first because $\beta < 1$ and the
second because $\beta > 0$. Its exponent is
\[
  F_{\mathrm I}(\sigma) = \sigma^\beta, \qquad \sigma \ge 0 .
\]
The profile is homogeneous: $k_{\mathrm I}(u/c) = c^\beta k_{\mathrm I}(u)$ for $c > 0$. So the
dilated datum $(c\,b_0, k_{\mathrm I}(\cdot/c))$ of the proof of \ref{prop:bridge-strictness}(4),
whose exponent is $F_{\mathrm I}(c\,\cdot)$, is the datum $(c^\beta b_0, c^\beta k_{\mathrm I})$,
whose exponent is $c^\beta F_{\mathrm I}$; hence $F_{\mathrm I}(\sigma) = \sigma^\beta
F_{\mathrm I}(1)$ for every $\sigma > 0$. The two sides agree on a neighbourhood of $\sigma = 1$,
so their derivatives there agree: the right side has derivative $\beta F_{\mathrm I}(1)$, and by
the derivative formula in the proof of \ref{prop:bridge-strictness}(2)
\[
  F_{\mathrm I}'(1) = \int_0^\infty e^{-u}\,c_\beta u^{-\beta}\,du = c_\beta\,\Gamma(1-\beta)
  = \beta .
\]
Hence $F_{\mathrm I}(1) = 1$, and $F_{\mathrm I}(0) = 0$ by the representation. By
\ref{prop:bridge-families}(4), $F_{\mathrm I}(\omega^2/2) = 2^{-\beta}|\omega|^{2\beta}$. Given
$0 < \alpha < 2$, take $\beta = \alpha/2$: then $|\omega|^\alpha = 2^{\alpha/2}F_{\mathrm
I}(\omega^2/2)$ is a positive multiple of a member of $\mathcal{S}$, and lies in $\mathcal{S}$,
which is a cone.
\end{proof}

% draft: Proposition 9.3(4), the membership sentence, its Student-t member (draft Lemma 9.3'')
% NEW as a node (2026-09-11, R121, split under the delegation). The sentence is not new: it was
% one of the four memberships of lem:subordinated-members (R44) and, before that, of
% prop:bridge-strictness(4). It is stated apart because the other three memberships are now
% machine-checked and this one cannot be on Lean core: the causal datum it needs is the \Levy{} data
% of the inverse-gamma delay law, which is what ledger A18 grounds, and A18 stays unadmitted (the
% author's decision of 2026-09-11: split, not admit). Typed as Skeleton.student_subordinated, the
% third conjunct of the former Skeleton.subordinated_members verbatim, which quantifies over the
% Student-t exponent through its transform.
% twin: hcs prop:bessel-family (causal ledger A8) for the causal datum.
% CHANGED (2026-09-15, module B step 2, R160): THE DELAY LAW BECOMES A HYPOTHESIS, and the node
% is [T] and \leanok. It read "S contains the Student-t family", [A] on ledger A18 and notready,
% waiting on Skeleton.student_causal -- the declaration the entry would have carried. What the
% proof consumes is not the named-class result (Halgreen: the generalized inverse Gaussian laws
% are generalized gamma convolutions) but one causally admissible datum whose delay law is the
% inverse-gamma law of shape a; quantifying over that datum loses nothing, and the passage from
% it to membership in S is the bridge image, machine-checked as SpatialLine.student_subordinated
% on Lean core. The existence of the datum is the cited remark rem:student-ggc, marked as not
% machine-checked, which is where A18 now binds; the entry stays unadmitted and the trust
% boundary is unchanged. Skeleton.student_subordinated and Skeleton.student_causal are withdrawn.
% The direction is the delegation's: the node claims less -- a conditional -- and what it claims
% is checked. (P-0010, which asked for exactly this to be tried before module B is drafted.)
\begin{lemma}[the Student-t family is subordinated]\label{lem:student-subordinated}
\uses{prop:bridge-strictness,prop:bridge-families,def:causal-admissible,prop:student-t}
\lean{SpatialLine.student_subordinated}\leanok
\notes{dickman-cin-rays}
Let $a > 0$, let $T_1$ be the inverse-gamma law of shape $a$ of \ref{prop:student-t}, and suppose
$T_1$ is the delay law of a causally admissible exponent: that for some $F_{\mathrm I}$ as in
\ref{def:causal-admissible},
\[
  \mathbb{E}e^{-\sigma T_1} = e^{-F_{\mathrm I}(\sigma)}, \qquad \sigma \ge 0 .
\]
Then $\mathcal{S}$ contains the Student-t family of shape $a$, with $F_{\mathrm I}$ as its causal
ancestor. That $T_1$ is such a delay law is cited at \ref{rem:student-ggc} and is not asserted
here.
\statusT\quad\emph{(Machine-checked to Lean core as
\texttt{SpatialLine.student\_subordinated}, at the reading the withdrawn skeleton target already
had: the Student-t exponent is quantified over through its transform --- every $F$ whose
exponential $e^{-F}$ is the cosine transform of the Student-t law of shape $a$ is subordinated
--- which loses nothing, the exponential being injective. The hypothesis is the whole of what
this node once took from ledger A18, and it is weaker than the entry: only self-decomposability
of $T_1$ is used here, not that it is a generalized gamma convolution, which
\ref{prop:student-t}(3) needs and \ref{rem:student-ggc} supplies. Nothing is admitted ---
\ledger{A18} grounds a hypothesis this node quantifies over, and no name of it is on the trust
boundary.
\emph{Provenance.} Split from \ref{lem:subordinated-members} on 2026-09-11 (R121), when that
node's other three memberships were machine-checked; restated with the delay law as a hypothesis
on 2026-09-15, under the safe-direction delegation.)}
\end{lemma}

\begin{proof}\leanok
Conditioning the Brownian kernel on $T_1$, the Student-t law of shape $a$ is the mixture of the
Gaussian laws against the law of $T_1$ --- \ref{prop:bridge-families}(3) --- so its cosine
transform at $\omega$ is $\mathbb{E}e^{-T_1\omega^2/2}$, which is the mixture identity of
\ref{lem:bridge-exponents}. The hypothesis read at $\sigma = \omega^2/2$ turns that into
$e^{-F_{\mathrm I}(\omega^2/2)}$. So an exponent $F$ of the Student-t family --- one with
$e^{-F(\omega)} = \hat\phi_1(\omega)$ at every frequency --- satisfies
$e^{-F(\omega)} = e^{-F_{\mathrm I}(\omega^2/2)}$ there, hence
$F(\omega) = F_{\mathrm I}(\omega^2/2)$ pointwise, the exponential being injective. That is
membership in $\mathcal{S}$, with $F_{\mathrm I}$ as the ancestor.
\end{proof}

% CHANGED 2026-09-19 (R178, second external review round, referee B finding 12): the paragraph
% below said "an admissible F whose CATALOGUE is carried by a compact subset of (0,infty)". The
% catalogue is the \Levy{} measure nu = k dx/x (rem:displacement-dictionary), and for a
% nonincreasing k not identically zero its support contains a neighbourhood of the origin, so no
% catalogue is carried by a compact subset of the open half-line and the sentence was empty as
% written. What the argument needs, and what it is true of, is the TAIL MEASURE varpi = -dk: a
% profile that is continuous and of bounded support. The conclusion is unchanged; a prose sentence
% is corrected.
The criterion in (3) is worth a second look, because it is the third appearance of one
phenomenon. A zero of an \emph{exponent} at $\omega_0 \ne 0$ means a lattice law
(\ref{lem:lattice-zero}) and is excluded by covariance (\ref{lem:no-lattice}); a zero of the
\emph{symbol} $B = \omega F'$ at $\omega_0 \ne 0$ means a lattice \emph{profile}, a catalogue
whose jump sizes lie on a lattice, and is excluded by no axiom --- the $\Cin$ rays are the case
of a single jump size. It is excluded exactly by subordination, because a Gaussian mixture
spreads every jump size over the whole line.

So the bridge is not onto, and the symmetric theory is strictly larger than the image of the
causal theory under it.
Symmetric self-decomposable laws with a lattice-step profile are members of the
spatial theory with no causal ancestor. That is the direction the statements above give --- it is
clause (3) in its general form, of which clause (4)'s exclusion of the $\Cin$ rays is the case of a
single jump size. The
converse --- that a member with no causal ancestor must have its catalogue mass on steps --- is
not asserted, and it fails: an admissible $F$ whose \emph{tail measure} $\varpi = -dk$ is carried
by a compact subset of $(0,\infty)$ and has no atom there --- equivalently, whose profile is
continuous and of bounded support --- has no ancestor either, because a subordinated
profile is a Gaussian mixture of the delay catalogue (\ref{eq:bridge-profile}) and is therefore
positive at every $x > 0$. (The catalogue itself, the \Levy{} measure $\nu = k\,dx/x$ of
\ref{rem:displacement-dictionary}, is carried by no compact subset of $(0,\infty)$: a
nonincreasing $k \not\equiv 0$ is positive on a neighbourhood of the origin. The sentence said
``catalogue'' until 2026-09-19; R178.) How much larger than the lattice-step members the complement of
$\mathcal{S}$ is, this chapter does not settle.

% NEW NODE (2026-09-17, the author's review of module B, R166): the dimension filtration as a
% statement. The author's finding at section 4 of module B: rem:dimension-filtration was "an
% enormous amount of technical prose where a large part reads as a proof", and that it is not
% machine-checked does not make it less of a proof; it must be a proposition with a proof, with
% the motivation before it. So the remark's proof-shaped paragraphs -- the classes, the nesting,
% the slice inside every dimension, what clause (3) says -- are this node and its proof, and the
% remark keeps what is a remark: the geometric readings, the two questions with the literature,
% the family behind the lemma and the numerics. It is an [A] node: the top clause rests on the
% radial theorem, which is the NEW ledger entry A23 (Schilling-Song-Vondracek Thm. 13.14 p. 212
% with the Bernstein representation Thm. 3.2 p. 21, both read by the librarian on 2026-09-17), and
% on A2 and A3 read at their anchors in dimension d; the slice clause rests on A3 in dimension d and
% on A7's polar criterion (Sato Thm. 15.10, already in A7's cite line). The subordination reading
% of the slice clause (Sato Thm. 30.1, (30.8), p. 198, read by the librarian the same day) is an
% interpretation and is not load-bearing: the proof exhibits the Levy pair by Tonelli, as
% lem:bridge-exponents does at d = 1. No Lean is scheduled; nothing here is machine-checked; the
% widening of the ledger (a new entry, A23) is the author's decision at the review.
% CHANGED (2026-09-18, module B register pass, R167): two hypotheses restored, both found by the
% blind reviewer of paper-b section 4. (a) prop:bridge-strictness(3)'s restating clause "that is,
% if and only if the profile k is a step function ..." dropped the conjunct a = 0 it restates
% (a omega^2 + Cin(tau omega) with a > 0 has a lattice-step profile and no stationary point; the
% Lean, SpatialLine.bridge_strictness_stationary, has P.a = 0 in the equivalence). (b) Clause (3)
% of this node excluded every admissible F with a lattice-step profile from A_infty, by a proof
% that goes through prop:bridge-strictness(3) and therefore needs a = 0; the clause, its summary
% sentence and the proof now carry "with Gaussian coefficient 0". (The claim for a > 0 is true by
% another route, the profile of a member of B being continuous, and is not asserted.) (c) The
% proof's use of the radial theorem said nothing of the dilation between F(|xi|) and
% f(|xi|^2) = F(sqrt 2 |xi|); one sentence now does. All three make the article claim the same or
% less; no status changes.
% CHANGED (2026-09-19, the author's decision D1 at the external review of module B, R172): A
% WIDENING. New clause (4), A_infty = S, with its proof. The article called this open
% (rem:dimension-filtration, the conclusions); the referee observed that it follows from the
% radial theorem applied to the REMAINDERS of self-decomposability, which are infinitely
% divisible in every dimension (Sato Prop. 15.5, read verbatim from the page image on
% 2026-09-19: "rho_b in (15.1) is uniquely determined and infinitely divisible", on R^d). The
% half-line step uses ledger A11 and A12 (a Bernstein function's exponential is completely
% monotone, hence a Laplace transform), A7 (Cor. 15.11) and Sato Thm. 24.11 with Remark 21.6
% (a law on the half-line: no Gaussian part, Levy measure on (0,oo), the Laplace form). The
% argument is the referee's; the checking and the writing are this session's.
% CHANGED 2026-09-19 (R179, second external review round, referee B finding 9): THREE STEPS OF THE
% PROOF ARE NAMED, none of them new. (a) Clause (4)'s "by the one-dimensional characterization" is
% Sato Cor. 15.11 p. 95 -- ledger A7's own anchor, for a law on the line -- and NOT
% prop:sd-exponents, which is stated for symmetric laws while the T of that step is carried by the
% half-line; the line now says so. (b) The radial theorem is stated for f : (0,infty) -> [0,infty),
% so g_c >= 0 is a hypothesis of its application: it is the monotonicity of the Bernstein function
% f at c^2 sigma <= sigma, half a sentence. (c) Sato's Def. 15.1 quantifies over b > 1 while the
% proof of clause (1) uses b in (0,1); the reciprocal is noted at the first use. No statement, no
% clause and no ledger line changes.
\begin{proposition}[the dimension filtration]\label{prop:dimension-filtration}
\uses{prop:bridge-strictness,def:causal-admissible,lem:bridge-exponents,prop:sd-exponents,lem:profile-integrability,prop:fourier-toolbox,prop:choquet-cone}
For an integer $d \ge 1$ let $\mathcal{A}_d$ be the set of admissible $F$ for which
$\xi \mapsto F(|\xi|)$, $\xi \in \RR^d$, is the exponent of a self-decomposable law on $\RR^d$;
let $\mathcal{A}_\infty$ be the intersection of all the $\mathcal{A}_d$; and let $\mathcal{B}$ be the set of
admissible $F$ for which $\sigma \mapsto F(\sqrt{2\sigma})$ is a Bernstein function on
$(0,\infty)$.
\begin{enumerate}
\item \emph{(The bottom, and the nesting.)} $\mathcal{A}_1$ is the whole admissible cone, and
$\mathcal{A}_{d+1} \subseteq \mathcal{A}_d$ for every $d \ge 1$.
\item \emph{(The slice sits in every dimension.)} Let $(b_0,k_{\mathrm I})$ be causally
admissible with exponent $F_{\mathrm I}$. For every $d$ the function
$\xi \mapsto F_{\mathrm I}(|\xi|^2/2)$ on $\RR^d$ is the exponent of an isotropic infinitely
divisible law, with Gaussian covariance $b_0 I$ and with a \Levy{} measure of radial density
$f_d(r) = \int_0^\infty (2\pi u)^{-d/2}e^{-r^2/2u}\,k_{\mathrm I}(u)\,du/u$, whose polar profile
is
\begin{equation}\label{eq:radial-profile}
  |S^{d-1}|\,r^d f_d(r)
   = \frac{2}{\Gamma(d/2)}\int_0^\infty v^{d/2-1}e^{-v}\,
     k_{\mathrm I}\Bigl(\frac{r^2}{2v}\Bigr)\,dv, \qquad r > 0 ,
\end{equation}
nonincreasing in $r$; so that law is self-decomposable, and
$\mathcal{S} \subseteq \mathcal{A}_d$. At $d = 1$ the polar profile is the folded profile
\ref{eq:bridge-profile}.
\item \emph{(The top.)} $\mathcal{A}_\infty \subseteq \mathcal{B}$. Consequently every
$F \in \mathcal{A}_\infty$ with $F \not\equiv 0$ has $F' > 0$ on $(0,\infty)$, and no admissible
$F \not\equiv 0$ with Gaussian coefficient $0$ whose profile is a step function with its jumps on a lattice --- in particular
no $\Cin$ ray --- lies in $\mathcal{A}_\infty$.
\item \emph{(The intersection is the slice.)} $\mathcal{A}_\infty = \mathcal{S}$: an admissible
$F$ whose radial extension is the exponent of a self-decomposable law on $\RR^d$ for every $d$
is subordinated.
\end{enumerate}
So $\mathcal{S} = \mathcal{A}_\infty \subseteq \dots \subseteq \mathcal{A}_2 \subseteq
\mathcal{A}_1$, and the lattice-step members of the line's theory with no Gaussian part lie outside
$\mathcal{A}_\infty$.
\statusA\quad\ledger{A23}\quad\ledger{A7}\quad\ledger{A3}\quad\ledger{A2}\quad\ledger{A11}\quad\ledger{A12}\quad\emph{(\textbf{Assignment.}
Clause (4), added on 2026-09-19, spends \ledger{A23} a second time, on the remainders of
self-decomposability; \ledger{A7} at the full letter of its Prop.~15.5, that the remainder
$\rho_b$ of a self-decomposable law on $\RR^d$ is infinitely divisible, and at its
one-dimensional characterization, read for a law carried by the half-line with the source's
description of such laws (no Gaussian part, \Levy{} measure on $(0,\infty)$, the Laplace form of
the exponent); \ledger{A2} in dimension $d$ for the remainders; and \ledger{A12} with
\ledger{A11}, that the exponential of the negative of a Bernstein function is completely
monotone and hence the Laplace transform of a positive measure. The argument of clause (4) is
the external referee's. For clauses (1) to (3):
\ledger{A23} carries the radial theorem at its source's letter --- $f : (0,\infty) \to [0,\infty)$
is a Bernstein function if and only if $\xi \mapsto f(|\xi|^2)$ is continuous and negative
definite on $\RR^d$ for all $d \in \NN$ --- and the representation of a Bernstein function,
$f(\sigma) = a + b\sigma + \int_{(0,\infty)}(1 - e^{-\sigma u})\,m(du)$ with a unique triplet;
clause (3) spends both. \ledger{A7} carries, beyond the line case the rest of the article uses,
the definition of a self-decomposable law on $\RR^d$, that such a law is infinitely divisible,
and the polar criterion, that an infinitely divisible law on $\RR^d$ is self-decomposable if and
only if its \Levy{} measure has the polar form with a nonincreasing radial profile along every
ray; clauses (1) and (2) spend it in that dimension. \ledger{A3} is read at its anchor in
dimension $d$, where the source states it: the \Levy--Khintchine representation on $\RR^d$, its
converse and its uniqueness; clause (2) spends the converse and the uniqueness. \ledger{A2} is
read at its anchor in dimension $d$ as well: the exponent of an infinitely divisible law on
$\RR^d$ is continuous and negative definite in the kernel sense, which is the hypothesis the
radial theorem asks for; clause (3) spends it. What is NOT carried, and not claimed: any
axiomatics in dimension $d$ --- the classes $\mathcal{A}_d$ are defined by the exponent
condition and not by families of operators, and the reading of them as the radial restrictions
of isotropic admissible families is \ref{rem:dimension-filtration}'s; the subordination reading
of clause (2), that the law is Brownian motion on $\RR^d$ run for the delay, which the proof
does not use; and anything about the top of the filtration beyond the inclusion, whether
$\mathcal{A}_\infty = \mathcal{S}$ being open. Nothing here is machine-checked and no Lean is
scheduled: the article has no development of laws on $\RR^d$.)}
\end{proposition}

\begin{proof}
(1) An admissible $F$ is of the form \ref{eq:sd-profile}, which is the symmetric
\Levy--Khintchine form of a self-decomposable law on the line (\ref{prop:sd-exponents}), so
$F \in \mathcal{A}_1$; and $\mathcal{A}_1$ consists of admissible $F$ by definition. For the
nesting let $F \in \mathcal{A}_{d+1}$, let $\mu$ be a self-decomposable law on $\RR^{d+1}$ with
exponent $\xi \mapsto F(|\xi|)$, and let $\Pi : \RR^{d+1} \to \RR^d$ be the projection that
drops the last coordinate, so that $\Pi^{\mathsf T}$ is an isometric embedding. The image law
$\Pi_*\mu$ has transform $\eta \mapsto \hat\mu(\Pi^{\mathsf T}\eta) = e^{-F(|\eta|)}$, so its
exponent is $\eta \mapsto F(|\eta|)$ on $\RR^d$. It is self-decomposable: for $b \in (0,1)$ the
definition gives a law $\rho_b$ on $\RR^{d+1}$ with $\hat\mu(\xi) = \hat\mu(b\xi)\,\hat\rho_b(\xi)$
--- the source states it at the reciprocal parameter, quantifying over $1/b > 1$ ---
for every $\xi$, and reading this at $\xi = \Pi^{\mathsf T}\eta$ gives
$\widehat{\Pi_*\mu}(\eta) = \widehat{\Pi_*\mu}(b\eta)\,\widehat{\Pi_*\rho_b}(\eta)$, which is the
definition on $\RR^d$ with $\Pi_*\rho_b$ as the residual. So $F \in \mathcal{A}_d$.

(2) Write $g^{(d)}_u$ for the centred Gaussian density on $\RR^d$ of covariance $u\,I$, and put
$\nu_d(dx) := \bigl(\int_0^\infty g^{(d)}_u(x)\,k_{\mathrm I}(u)\,du/u\bigr)\,dx$ on
$\RR^d \setminus \{0\}$, a positive measure with radial density $f_d$. It is a \Levy{} measure:
since $\int_{\RR^d}(1 \wedge |x|^2)\,g^{(d)}_u(x)\,dx$ is at most the total mass $1$ and at most
the trace $du$ of the covariance, Tonelli gives
$\int (1 \wedge |x|^2)\,\nu_d(dx) \le \int_0^\infty (1 \wedge du)\,k_{\mathrm I}(u)\,du/u
 \le d\int_0^\infty (1 \wedge u)\,k_{\mathrm I}(u)\,du/u < \infty$, the last integral being the
single-integral reading of the two conditions of \ref{def:causal-admissible}. Since
$e^{-u|\xi|^2/2} = \int_{\RR^d}\cos(\xi\cdot x)\,g^{(d)}_u(x)\,dx$ and $\int g^{(d)}_u = 1$, a second
Tonelli exchange gives
\[
  F_{\mathrm I}\bigl(\tfrac12|\xi|^2\bigr)
   = \tfrac12\,\xi\cdot(b_0 I)\xi
     + \int_{\RR^d}\bigl(1 - \cos\xi\cdot x\bigr)\,\nu_d(dx), \qquad \xi \in \RR^d,
\]
the symmetric \Levy--Khintchine form on $\RR^d$ with Gaussian covariance $b_0 I$ and \Levy{}
measure $\nu_d$; both are invariant under the orthogonal group. By the converse clause of the
representation theorem in dimension $d$ this is the exponent of an infinitely divisible law,
isotropic because its transform is, and by the uniqueness clause that law's \Levy{} measure is
$\nu_d$. Its polar form is read off the radial density: in polar coordinates
$\nu_d(dx) = f_d(r)\,r^{d-1}\,dr\,d\sigma(\theta)$, so the profile along every ray, in the
sense of the polar criterion, is $r \mapsto r^d f_d(r)$ up to the constant $|S^{d-1}|$. The
substitution $u = r^2/2v$, which leaves $du/u$ invariant, turns
$|S^{d-1}|\,r^d\int_0^\infty (2\pi u)^{-d/2}e^{-r^2/2u}\,k_{\mathrm I}(u)\,du/u$ into
\ref{eq:radial-profile}, using $|S^{d-1}| = 2\pi^{d/2}/\Gamma(d/2)$ and
$(2\pi u)^{-d/2}r^d = \pi^{-d/2}v^{d/2}$. The weight $v^{d/2-1}e^{-v}$ no longer involves $r$,
and $k_{\mathrm I}(r^2/2v)$ is nonincreasing in $r$ for each fixed $v$ because $k_{\mathrm I}$ is
nonincreasing; so the polar profile is nonincreasing, and by the polar criterion the law is
self-decomposable. Its exponent restricted to a line is $\omega \mapsto F_{\mathrm I}(\omega^2/2)$,
which is admissible by \ref{lem:bridge-exponents}, so every subordinated $F$ lies in
$\mathcal{A}_d$. At $d = 1$, $|S^0| = 2$ and \ref{eq:radial-profile} is \ref{eq:bridge-scaling}
with the folding factor in place, that is the folded profile \ref{eq:bridge-profile}.

(3) Let $F \in \mathcal{A}_\infty$ and fix $d$. The law on $\RR^d$ with exponent
$\xi \mapsto F(|\xi|)$ is self-decomposable, hence infinitely divisible, so $e^{-tF(|\xi|)}$ is
a characteristic function for every $t > 0$, hence positive definite, and $\xi \mapsto F(|\xi|)$
is continuous and negative definite in the kernel sense. That holds for all $d$, and negative
definiteness is unchanged by the dilation $\xi \mapsto \sqrt2\,\xi$, so with
$f(\sigma) := F(\sqrt{2\sigma})$ the function $\xi \mapsto f(|\xi|^2) = F(\sqrt2\,|\xi|)$ is
continuous and negative definite on every $\RR^d$. By the radial
theorem $f$ is a Bernstein function on $(0,\infty)$, and
$F \in \mathcal{B}$. For the derivative, the representation gives
$f(\sigma) = a + b\sigma + \int_{(0,\infty)}(1 - e^{-\sigma u})\,m(du)$ with $a, b \ge 0$ and
$\int (1 \wedge u)\,m(du) < \infty$; $a = \lim_{\sigma \downarrow 0} f(\sigma) = F(0) = 0$, $F$
being continuous with $F(0) = 0$. On $[\sigma_0,\infty)$ for any $\sigma_0 > 0$ the
$\sigma$-derivative $u\,e^{-\sigma u}$ of the integrand is dominated by $u$ on $(0,1]$ and by
$1/(e\sigma_0)$ on $(1,\infty)$, both $m$-integrable, so
$f'(\sigma) = b + \int_{(0,\infty)} u\,e^{-\sigma u}\,m(du)$ for $\sigma > 0$, a sum of
nonnegative terms that vanishes only if $b = 0$ and $m = 0$, that is only if $f \equiv 0$, that is
only if $F \equiv 0$. So $F \not\equiv 0$ gives $f' > 0$, and the chain rule gives
$F'(\omega) = \omega\,f'(\omega^2/2) > 0$ for $\omega > 0$. Finally, an admissible $F \not\equiv 0$
with Gaussian coefficient $0$ whose profile is a step function with its jumps on a lattice has, by
\ref{prop:bridge-strictness}(3), a stationary point $\omega_0 \ne 0$, and by evenness one with
$\omega_0 > 0$; that contradicts $F' > 0$, so no such $F$ lies in $\mathcal{A}_\infty$. The
$\Cin$ rays are the case of a single jump.

(4) $\mathcal{S} \subseteq \mathcal{A}_\infty$ is clause (2). Conversely let
$F \in \mathcal{A}_\infty$ and put $f(\sigma) := F(\sqrt{2\sigma})$, a Bernstein function with
$f(0{+}) = 0$ by clause (3) and the continuity of $F$. Fix $0 < c < 1$. For each $d$ let $\mu_d$
be the self-decomposable law on $\RR^d$ with exponent $\xi \mapsto F(|\xi|)$. By the definition
there is a law $\rho_{c,d}$ with $\hat\mu_d(\xi) = \hat\mu_d(c\xi)\,\hat\rho_{c,d}(\xi)$, and by
the cited proposition $\rho_{c,d}$ is infinitely divisible. Its transform is
$e^{-(F(|\xi|) - F(c|\xi|))}$, the transform of $\mu_d$ having no zero, so its exponent, being
continuous with value $0$ at the origin, is $\xi \mapsto F(|\xi|) - F(c|\xi|)$. That function
is therefore continuous and negative definite on $\RR^d$, for every $d$, and so is its dilate
$\xi \mapsto F(\sqrt2|\xi|) - F(c\sqrt2|\xi|) = g_c(|\xi|^2)$ with
$g_c(\sigma) := f(\sigma) - f(c^2\sigma)$. That $g_c$ takes values in $[0,\infty)$, which the
radial theorem asks of its argument, is the monotonicity of the Bernstein function $f$ at
$c^2\sigma \le \sigma$. By the radial theorem, applied as in clause (3),
$g_c$ is a Bernstein function, and $g_c(0{+}) = 0$.

The exponential of the negative of a Bernstein function is completely monotone, and a completely
monotone function is the Laplace transform of a positive measure on $[0,\infty)$, whose mass is
the value at $0{+}$. So there are probability laws $T$ and $R_c$ on $[0,\infty)$ with Laplace
transforms $e^{-f}$ and $e^{-g_c}$, and $e^{-f(\sigma)} = e^{-f(c^2\sigma)}\,e^{-g_c(\sigma)}$
says that $T$ has the law of $c^2T' + R$ with $T'$ distributed as $T$, $R$ as $R_c$, and the two
independent, a Laplace transform determining a law on the half-line. As $c$ ranges over $(0,1)$
so does $c^2$, so $T$ is self-decomposable. By the one-dimensional characterization of
self-decomposable \Levy{} measures --- A7 at its own anchor, \sscite{sato1999levy},
Corollary~15.11, p.~95, which is the two-sided form for a law on the line and not
\ref{prop:sd-exponents}, this $T$ being one-sided rather than symmetric --- its \Levy{}
measure is $k_{\mathrm I}(u)\,du/u$ on $(0,\infty)$ with $k_{\mathrm I}$ nonincreasing, the law
being carried by the half-line, and its Laplace exponent is
$f(\sigma) = b_0\sigma + \int_0^\infty(1 - e^{-\sigma u})\,k_{\mathrm I}(u)\,du/u$ with
$b_0 \ge 0$ and $\int_0^\infty(1 \wedge u)\,k_{\mathrm I}(u)\,du/u < \infty$, which is the
single-integral reading of the two conditions of \ref{def:causal-admissible}. So $f$ is causally
admissible, and $F \in \mathcal{S}$ by \ref{prop:bridge-strictness}(1).
\end{proof}

% NEW as a remark (2026-09-15, module B step 2, R161), from the postdoc's proposal P-0009 and
% written for module B's literature paragraph. It is a remark and not a statement node on purpose:
% every sentence below is either cited to a held page or an elementary consequence of nodes already
% in this chapter, and the two questions the proposal raises are left open in the text. Nothing here
% is formalised and nothing here is admitted. The RADIAL form of Schoenberg's theorem is in no
% ledger entry -- A1 is Bochner (Sato pp. 8-9), A2 is the kernel-versus-exponential form (Schilling
% p. 36) -- and the closing paragraph says what an entry would have to carry if the reading ever
% became a statement node.
% FOUND WHILE WRITING IT, AND DELIBERATELY NOT ASSERTED AS A CLAIM OF THE CHAPTER: an explicit
% admissible exponent of type G whose mixing subordinator is not self-decomposable, which answers
% question (i) in the negative, and the same one-parameter family witnessing strictness of the
% filtration at d = 2. Both are in the closing paragraph, in full in Formalization/SKELETON.md
% section 23, and they are candidate statement nodes for the author.
% CHANGED (2026-09-15, module B step 2, R161): this remark is new; no statement of this chapter is
% touched, and no status changes.
% CHANGED (2026-09-15, module B step 2, R163): the two constructions ARE NOW NODES, on the author's
% decision of the same day, and two paragraphs are rewritten to cite them. The "two questions"
% paragraph no longer says that no statement of this chapter answers either: lem:filtration-strictness
% answers both, the first in the negative and the second at the bottom of the filtration, and what
% is left open is A_infty = S, whose argument is a sketch (the convolution reading of
% eq:radial-profile and an approximate identity in the dimension) and is stated here as a sketch.
% The "two constructions, offered and not asserted" paragraph becomes the family and the numerics:
% the general (beta, a, b) member and the thresholds c_d stay here, since they are numerical, and
% the memberships at c = 4 and c = 8 move to the lemma, where each is a checkable inequality. The
% comparison of the two sufficient conditions at d = 1 is new and is why the lemma's hypothesis is
% c <= 6 + 4 sqrt 2 rather than the remark's c <= e/2: the sharper bound keeps the negative half of
% the derivative and covers both values of c, so clause (2) inherits admissibility from clause (1)
% instead of repeating it. No status of any node changes; nothing here is machine-checked.
% CHANGED (2026-09-17, the author's review of module B, R166): the proof-shaped paragraphs -- the
% radial theorem as applied, the classes, the nesting, why those classes are the restrictions,
% the slice sits inside every dimension, what clause (3) says, the status paragraph -- are now
% prop:dimension-filtration above and its proof, [A] on the new ledger entry A23 with A2, A3 and A7
% read in dimension d. What stays is what is a remark: the geometric reading (the type G
% identification and the radial-restriction heuristic, now marked as a reading), the two questions
% with the literature, and the family behind the lemma with its numerics. No status changes; the
% annotations of def:delay-data and lem:filtration-strictness now cite the proposition for what is
% asserted and the remark for what stays a reading.
% CHANGED (2026-09-19, article mirror of module B, R174): the geometric reading's pushforward
% argument was one sentence of about 120 words and is now five (the register pass's item 20), and
% its closing clause no longer says twice that the reading is a reading -- the paragraph is headed
% "The geometric reading" and opens with "heuristically". No claim changes; nothing is added or
% dropped.
\begin{remark}[the dimension filtration: readings, questions, numerics]\label{rem:dimension-filtration}
\uses{prop:dimension-filtration,prop:bridge-strictness,def:causal-admissible,lem:bridge-exponents}
\emph{The geometric reading.} By the radial theorem, $\mathcal{B}$ is exactly the set of
admissible $F$ whose radial extension $\xi \mapsto F(|\xi|)$ is a \Levy{} exponent on $\RR^d$
for every $d$ at once, and the laws so obtained are the variance mixtures of the centred normal
law with an infinitely divisible mixing law, the type G laws, the mixing law being the delay of
this chapter (\sscite{maejima2002type}, (1.1), p.~323, and for $\RR^d$ p.~325, who attribute the class to
Marcus and to Rosi\'nski; \sscite{aoyama2007characterizations}, Definition~1.1, p.~148; the
one-dimensional mixture form is also \sscite{steutel2004infinite}, (2.7), p.~345); the criterion is about all
dimensions together, the hard direction of the theorem taking a limit in the dimension. The
classes $\mathcal{A}_d$ read, heuristically, as the radial restrictions of the isotropic
admissible families of $\RR^d$. Take a family of operators on $L^1(\RR^d)$ satisfying the evident
analogue of (A1)--(A8) and (ND), with kernels invariant under the orthogonal group, and push the
kernels forward under an orthogonal projection onto a line. What comes out satisfies (A1)--(A8)
and (ND) on the line, with the restriction of $\xi \mapsto F(|\xi|)$ as its exponent: pushforward
under a linear map commutes with convolution and carries probability measures to probability
measures and dilations to dilations; the reflection of the line is induced by an orthogonal map;
and a rotation invariant measure carried by a hyperplane is $\delta_0$. A linear
image of a self-decomposable law is self-decomposable, which is the nesting argument of
\ref{prop:dimension-filtration}(1). The converse would need the analogue of
\ref{thm:main-characterization} in dimension $d$, which this article does not prove. So the
classes are defined by the exponent condition and not by the families, and only the exponent
condition is used. In these terms \ref{prop:dimension-filtration}(3) says that the lattice-step members
of the line's theory, the $\Cin$ rays among them, are not radial restrictions of isotropic
admissible families valid in every dimension, and are outside the type G laws, not merely
outside $\mathcal{S}$.

\emph{Two questions, and what the lemma below settles.} The first is settled in the literature, by a theorem of Sato's; the second, to the author's knowledge, is not. The first is whether $\mathcal{S}$ is all of the type G slice, that is whether the
mixing subordinator of a self-decomposable Gaussian variance mixture is itself
self-decomposable. The held sources prove the implication in the direction this chapter uses
and not its converse. \sscite{halgreen1979self} announces on p. 14 and proves on pp. 15--16
that for a pure variance mixture self-decomposability of the mixing law implies
self-decomposability of the mixture, and the question left open there is a different one, the
mean-variance case. \sscite{barndorffnielsen1977infinite} gives infinite divisibility of the
mixture from infinite divisibility of the mixing law, in $\RR^r$ for every $r$ (pp. 309--311).
One level up the corresponding statement is an equivalence: $\varphi(s^2/2)$ is the transform
of a symmetric extended generalized gamma convolution if and only if $\varphi$ is the transform
of a generalized gamma convolution (\sscite{bondesson1992generalized}, Thm. 7.3.1, pp. 115--116,
proved in both directions), which is the Thorin-level twin of \ref{prop:thorin-subclass}(3). At the level of class L the converse is false, and this is known: \sscite{sato2001subordination}, Theorem~1.2, constructs a self-decomposable law of type G that is the Gaussian mixture of no self-decomposable law on the half-line, the mixing law being determined by the mixture (his Lemma~3.1).
% CHANGED 2026-09-19 (R192, left by batch S2; referee B finding 10(iii), with Sato's (3.7) read
% from the held copy this session): "what this paper adds is an explicit family with its
% constants" is withdrawn. Sato's own witness is explicit, and it is printed here; what is new is
% the crossing of the dimensions and the identification at the top. Safe direction: a novelty
% claim is dropped.
His witness is explicit as well: the delay profile of his (3.7), p.~8, is $2s^{-1/2}$ on
$(0,1/(a+b)) \cup [1/a,\infty)$ and $-2s^{-1/2} + c$ on $[1/(a+b),1/a)$, for fixed $a > b > 0$ and
a constant $c$ chosen to keep it nonnegative --- a power with a rising piece inserted, which is
the design of the family below with a slope where that family has a jump. What is new here is
therefore not the explicitness. It is the crossing of the dimensions one at a time,
\ref{lem:filtration-strictness}(3), and the identification at the top,
\ref{prop:dimension-filtration}(4):
\ref{lem:filtration-strictness}(1) exhibits an admissible exponent that is of the type G form of
\ref{def:delay-data} and is not subordinated. Reading that as an answer to the question as
posed uses the radial theorem, which makes the type G form the class of variance mixtures, and
the existence of the delay law at \ref{rem:bridge-subordination}, which makes the mixing law a
law at all. The second question is whether the filtration is strict inside the admissible cone,
and in particular whether $\mathcal{A}_2 = \mathcal{A}_\infty$. It is strict at the bottom, which
is \ref{lem:filtration-strictness}(3): one admissible exponent lies in $\mathcal{A}_2$ and not in
$\mathcal{A}_3$, another in $\mathcal{A}_1$ and not in $\mathcal{A}_2$, and the first separates
$\mathcal{A}_2$ from $\mathcal{A}_\infty$ as well. Admissibility in the
plane therefore does not already force the subordination picture. The top of the filtration is settled by \ref{prop:dimension-filtration}(4):
$\mathcal{A}_\infty = \mathcal{S}$, by the radial theorem applied to the remainders of
self-decomposability.
% CHANGED 2026-09-19 (R193, left by batch S2; referee E finding 4.2 asked for the sentence, and
% the librarian's search of the same day is what it reports: notes/reviews/cone/
% SOURCES-round2-2026-09-19.md, Part 3).
% 2026-09-20: the TODO left here for the librarian is done. The six papers were acquired by the
% author and read in full (notes/reviews/cone/SOURCES-round2-new-papers-2026-09-20.md); the
% sentence below reports that reading.
\emph{We have not found this equivalence stated.} The sufficiency half is classical and is in
print more than once: in dimension one (\sscite{halgreen1979self}; \sscite{sato2001subordination},
Theorem~1.1); for normal variance mixtures in a dimension $r$ that does not enter the hypothesis
on the mixing law, which are self-decomposable when the mixing law is, and ``self-decomposable in
any dimension'' for the generalized hyperbolic family at zero drift
(\sscite{barndorffnielsen1982normal}, pp.~147 and~149); and in a fixed dimension for an
operator-stable subordinand (\sscite{barndorffnielsen2001multivariate}, Theorem~6.1, whose
Example~6.2 shows that self-decomposability of the subordinand alone does not suffice). None of
these states a converse. The literature on type G laws on $\RR^d$
(\sscite{maejima2002type}; \sscite{aoyama2007characterizations};
\sscite{barndorffnielsen2006some}), read in full, takes an infinitely divisible mixing law, fixes the
dimension, and does not treat the question. So the direction from self-decomposability in every
dimension to a self-decomposable mixing law is, as far as these sources go, stated here for the
first time; the sentence records what was read, on 2026-09-19 and 2026-09-20, and not an
exhaustive search. In the coordinate $x = \log r^2$ the same fact has a second reading:
\ref{eq:radial-profile} says that the polar profile is the convolution of the delay measure with
the law of $\log w$, $w$ a chi-square variable with $d$ degrees of freedom, and those laws
recentred at $\log d$ are an approximate identity, which is a second way to see why monotonicity in every dimension forces monotonicity of $k_{\mathrm I}$. That route is not carried out here, and clause (4) does not need it.

\emph{The family behind the lemma, and the numerics.} Writing the reading above produced an
answer to the first question and a witness for the second, and both are
\ref{lem:filtration-strictness}, stated at one parameter point of the following family. Take
$0 < \beta < 1$, $0 < a < b$, $c > 0$ and the delay profile
$k_{\mathrm I}(u) = u^{-\beta} + c\,\mathbf{1}_{[a,b]}(u)$. It satisfies the two integrability
conditions of \ref{def:causal-admissible} and has no nonincreasing version, so its
$F_{\mathrm I}$ is a Bernstein function and is not causally admissible. By
\ref{eq:radial-profile} the polar profiles of the associated isotropic laws are
\[
  2^{1+\beta}\,\frac{\Gamma(d/2+\beta)}{\Gamma(d/2)}\,r^{-2\beta}
   + 2c\Bigl[P\Bigl(\frac d2,\frac{r^2}{2a}\Bigr)
            - P\Bigl(\frac d2,\frac{r^2}{2b}\Bigr)\Bigr],
\]
$P$ here being the regularized lower incomplete gamma function. The lemma takes
$\beta = \tfrac12$, $a = 1$, $b = 2$, where the bracket is
$\operatorname{erf}(r/\sqrt2) - \operatorname{erf}(r/2)$ at $d = 1$ and
$e^{-r^2/4} - e^{-r^2/2}$ at $d = 2$.

What stays here rather than in the lemma is the numerics. At those parameters the polar profile
is nonincreasing in dimension $d$ exactly for $c \le c_d$, with $c_1 \approx 12.61$,
$c_2 \approx 5.24$, $c_3 \approx 3.21$ and $c_d \to 0$; the thresholds were computed numerically
from those one-variable inequalities, and the lemma asserts memberships only at $c = 4$ and
$c = 8$, where each inequality is elementary and is proved. Two sufficient conditions at $d = 1$
are worth comparing, since the second is what lets one hypothesis carry both values of $c$.
Dropping the negative half of the derivative of the bump term leaves $c\,r^2e^{-r^2/2} \le 1$ for
every $r$, that is $c \le e/2$; keeping it and bounding $\log(1/p) \le (1-p)/\sqrt p$ leaves
$c(\tfrac32 - \sqrt2) \le 1$, that is $c \le 6 + 4\sqrt2 \approx 11.66$, which is within eight per
cent of $c_1$ and is the hypothesis of \ref{lem:filtration-strictness}(1). Because $c_d \to 0$ the
family leaves every $\mathcal{A}_d$ eventually while staying of type G for every $c > 0$, which is
why $\mathcal{A}_\infty$ is not the type G slice.
\end{remark}

% NEW NODE (2026-09-15, module B step 2, R163): the vocabulary the two constructions need in order
% to be statements. It is def:causal-admissible with the monotonicity of the profile dropped and the
% measurability that field supplied put in its place, together with the polar profile the remark
% above computes. Nothing cited enters it: the two readings that make the classes geometric -- the
% radial theorem and Sato's polar criterion -- are the remark's, are cited there, and are not
% asserted here, so this node opens no ledger entry and needs none. The Lean type is
% Skeleton.DelayDatum; it carries no sorry, a structure being a definition, so the node is \leanok.
% Placed after the remark because the remark is where the reading lives; the definition is used
% only by lem:filtration-strictness below.
\begin{definition}[delay data, the type G exponents, and the polar profiles]
\label{def:delay-data}
\uses{def:causal-admissible}
\lean{Skeleton.DelayDatum, Skeleton.DelayDatum.exponent, Skeleton.IsTypeG,
Skeleton.polarProfile}\leanok
A \emph{delay datum} is a pair $(b_0, k_{\mathrm I})$ with $b_0 \ge 0$ and
$k_{\mathrm I} : (0,\infty) \to [0,\infty)$ measurable, subject to
$\int_0^1 k_{\mathrm I}(u)\,du < \infty$ and $\int_1^\infty k_{\mathrm I}(u)\,u^{-1}du < \infty$.
Its exponent is
\[
  F_{\mathrm I}(\sigma) = b_0\,\sigma
   + \int_0^\infty \bigl(1 - e^{-\sigma u}\bigr)\,k_{\mathrm I}(u)\,\frac{du}{u},
  \qquad \sigma \ge 0 ,
\]
and its \emph{polar profile in dimension $d$}, for an integer $d \ge 1$, is
\begin{equation}\label{eq:polar-profile}
  k_d(r) := \frac{2}{\Gamma(d/2)}\int_0^\infty v^{d/2-1}e^{-v}\,
   k_{\mathrm I}\Bigl(\frac{r^2}{2v}\Bigr)\,dv, \qquad r > 0 .
\end{equation}
Write $\mathcal{G}$ for the set of functions of the form $F(\omega) = F_{\mathrm I}(\omega^2/2)$,
$\omega \in \RR$, for a delay datum. These are type G exponents, the exponents of Gaussian
variance mixtures with an infinitely divisible mixing law, and they are those whose mixing law
has a \Levy{} measure of the form $k_{\mathrm I}(u)\,du/u$; a mixing law with an atomic \Levy{}
measure gives a type G exponent that is not in $\mathcal{G}$.
\statusT\quad\emph{(Vocabulary, and no claim. This is \ref{def:causal-admissible} with the
monotonicity of the delay profile dropped and the measurability that field supplied put in its
place, so every causally admissible datum is a delay datum and
$\mathcal{S} \subseteq \mathcal{G}$. Two geometric readings are not asserted here: that
$\mathcal{G}$ is the class of exponents of the symmetric type G laws, the variance mixtures of
the centred normal law with an infinitely divisible mixing law, which is the radial theorem
(ledger A23, spent by \ref{prop:dimension-filtration}(3)) together with the existence of the
delay law at \ref{rem:bridge-subordination}, and is \ref{rem:dimension-filtration}'s reading;
and that $k_d$ being nonincreasing is self-decomposability on $\RR^d$ of the isotropic law with
radial exponent $F$, which is the polar criterion, spent by \ref{prop:dimension-filtration}(2)
for a causally admissible datum and not asserted here for a general one. The definition reads
only the displayed forms, so no dimension theory and no ledger entry enters it. \ref{eq:polar-profile} is \ref{eq:radial-profile}, restated so that the lemma
below stands on its own; at $d = 1$ it is \ref{eq:bridge-scaling} with the folding factor
$2 = |S^0|$ in place, so $k_1$ is the folded spatial profile \ref{eq:bridge-profile} of $F$. The
Lean type is \texttt{Skeleton.DelayDatum}, which is \texttt{CausalAdmissible} without its
\texttt{k\_antitone} field, and \texttt{Skeleton.DelayDatum.ofCausal} is the inclusion.)}
\end{definition}

% NEW NODE (2026-09-15, module B step 2, R163): the two constructions of the P-0009 pass, stated on
% the author's decision of 2026-09-15. Typed as Skeleton.filtration_strictness_typeG and
% Skeleton.filtration_strictness_dimension, both sorry, so the node is \notready and the proof below
% is the proof of record. It spends NO ledger entry: what the argument needs where a reader would
% expect Bernstein's theorem is uniqueness of the DELAY DATA, and prop:bridge-strictness(1) does not
% carry it (it turns membership in S into an existential over data) while ledger A11 carries it only
% at its citation -- the admitted statement there is the existence equivalence alone. This article
% proves what is needed, as prop:laplace-uniqueness-locally-finite, which is [T] and Lean core; the
% derivative formula and the dominated limit at sigma -> infinity are the rest. Finding recorded in
% Formalization/SKELETON.md section 27.
% The hypothesis of clause (1) is c <= 6 + 4 sqrt 2 rather than the remark's c <= e/2 because the
% sharper elementary bound is what covers c = 4 and c = 8, so clause (2) does not have to repeat the
% admissibility argument. See the remark for the comparison and for the numerical thresholds.
% CHANGED 2026-09-19 (R191; the author's decision D3, a widening authorized at the second external
% review round of module B): CLAUSE (3) IS ADDED, and with it the node becomes [A]. The strictness
% of the chain A_1 > A_2 > A_3 was drawn in rem:dimension-filtration, in a figure caption and in
% the article's section 1, and was a clause of no statement (referee B finding 5). It is one now,
% and the referee's two missing steps are written into the proof: that k_d^{(c)} IS the polar
% profile of the law with radial exponent F^{(c)}(|xi|) -- the d-dimensional computation of
% prop:dimension-filtration(2) reads monotonicity only in its last line, which is now said -- and
% that a rotation-invariant \Levy{} measure of a self-decomposable law has a nonincreasing radial
% profile, which is Sato's polar criterion integrated over the sphere. Clause (3) therefore spends
% A7 (the criterion, both directions) and A3 (the representation in dimension d, converse and
% uniqueness), exactly the entries prop:dimension-filtration spends, and the status line changes
% from \statusT to \statusA. NO LEAN CHANGES: the two sorry'd declarations state clauses (1) and
% (2), clause (3) is prose-only, and the annotation says so. The \uses edge to
% prop:dimension-filtration is added, the classes A_d being defined there.
\begin{lemma}[the slice inside the type G exponents, and the polar profiles across dimensions]
\label{lem:filtration-strictness}
\uses{def:delay-data,def:causal-admissible,lem:bridge-exponents,lem:profile-integrability,prop:bridge-strictness,prop:laplace-uniqueness-locally-finite,thm:main-characterization,prop:dimension-filtration}
\lean{Skeleton.filtration_strictness_typeG, Skeleton.filtration_strictness_dimension}\notready
For $c > 0$ let $D^{(c)}$ be the delay datum with drift $0$ and profile
\[
  k^{(c)}_{\mathrm I}(u) = u^{-1/2} + c\,\mathbf{1}_{[1,2]}(u), \qquad u > 0 ,
\]
let $F^{(c)}_{\mathrm I}$ be its exponent, put
$F^{(c)}(\omega) := F^{(c)}_{\mathrm I}(\omega^2/2)$, and write $k^{(c)}_d$ for its polar profile
\ref{eq:polar-profile}.
\begin{enumerate}
\item \emph{The subordinated slice is strictly smaller than the type G exponents.}
$\mathcal{S} \subseteq \mathcal{G}$; and for $0 < c \le 6 + 4\sqrt2$ the function $F^{(c)}$ is
admissible in the sense of \ref{thm:main-characterization}, with Gaussian coefficient $0$ and
folded profile $k^{(c)}_1$, and lies in $\mathcal{G} \setminus \mathcal{S}$. Hence
$\mathcal{S} \subsetneq \mathcal{G}$.
\item \emph{The polar profiles cross the dimensions one at a time.} For $c = 4$ the profiles
$k^{(4)}_1$ and $k^{(4)}_2$ are nonincreasing on $(0,\infty)$ and $k^{(4)}_3$ is not; for $c = 8$
the profile $k^{(8)}_1$ is nonincreasing on $(0,\infty)$ and $k^{(8)}_2$ is not.
\item \emph{The filtration is strict at the bottom.} In the notation of
\ref{prop:dimension-filtration}, $F^{(8)} \in \mathcal{A}_1 \setminus \mathcal{A}_2$ and
$F^{(4)} \in \mathcal{A}_2 \setminus \mathcal{A}_3$. Hence
$\mathcal{A}_1 \supsetneq \mathcal{A}_2 \supsetneq \mathcal{A}_3$, and $F^{(4)}$ separates
$\mathcal{A}_2$ from $\mathcal{A}_\infty$ as well, that class lying inside $\mathcal{A}_3$.
\end{enumerate}
\statusA\quad\ledger{A7}\quad\ledger{A3}\quad\emph{(\textbf{Assignment.} Clauses (1) and (2) are
$[T]$ and spend no ledger entry; clause (3), added on 2026-09-19 (R191, the author's decision D3),
reads clause (2) through two cited facts in dimension $d$, both of them the ones
\ref{prop:dimension-filtration} spends. \ledger{A7} carries the definition of a self-decomposable
law on $\RR^d$ and the polar criterion --- that an infinitely divisible law on $\RR^d$ is
self-decomposable if and only if its \Levy{} measure has the polar form with a nonincreasing
radial profile along every ray --- at both of its directions: the sufficiency for the two
memberships, the necessity for the two exclusions. \ledger{A3} is read at its anchor in dimension
$d$, at the converse clause of the \Levy--Khintchine representation, which turns the computed
triplet into a law, and at its uniqueness clause, which identifies the \Levy{} measure of any law
with the given exponent. What they do not carry, and what the proof below writes out, are the two
steps the external review found missing: that $k^{(c)}_d$ \emph{is} the polar profile of the law
with radial exponent $F^{(c)}(|\xi|)$, which is the computation of
\ref{prop:dimension-filtration}(2) read at a datum with no monotonicity --- that computation uses
monotonicity only in its last line --- and that a rotation-invariant \Levy{} measure of a
self-decomposable law has a nonincreasing radial profile, which is the polar criterion averaged
over the sphere. Both are $[T]$.
\textbf{Clause (3) is prose-only}: the two typed declarations state clauses (1) and (2), and
nothing of clause (3) is typed, this article having no theory of self-decomposable laws on
$\RR^d$ --- which is also why clauses (1) and (2) are stated at the polar profiles, the objects
the criterion evaluates, and clause (3) at the classes.
\textbf{Stated, not machine-checked.} Both earlier clauses are typed, as
\texttt{Skeleton.filtration\_strictness\_typeG} and
\texttt{Skeleton.filtration\_strictness\_dimension}, and both carry \texttt{sorry}; the proof
below is the proof of record. \textbf{Clauses (1) and (2) spend no ledger entry.} Where a reader
expects Bernstein's theorem --- to rule out every causally admissible datum with the given
exponent, clause (1) of \ref{prop:bridge-strictness} being an existential over data --- what is
needed is uniqueness of the delay data, and A11 is admitted at its existence equivalence only. The
uniqueness this article proves, \ref{prop:laplace-uniqueness-locally-finite}, supplies it, with
the derivative formula from the proof of \ref{prop:bridge-strictness}(2) and one dominated limit.
\textbf{One geometric reading is still not asserted here}: that $\mathcal{G}$ is the class of type
G exponents in the sense of variance mixtures is the radial theorem (A23), and it is
\ref{rem:dimension-filtration}'s reading, not a clause of this lemma. The other one is clause (3)
since 2026-09-19. \textbf{The hypothesis of
clause (1)} is the exact threshold of the elementary bound in the proof, and is what carries both
$c = 4$ and $c = 8$; the cruder condition $c \le e/2$ of \ref{rem:dimension-filtration}, which
drops the negative half of one derivative, would not, and the numerical threshold, close to
$12.61$, is recorded there. \textbf{Why \texttt{notready}, and the price} (2026-09-15, written at
the statement): about $550$--$850$ Lean lines, in six pieces --- the datum, the closed form of the
polar profile, the admissibility clause, the monotonicity bounds, the two explicit failures, and
the uniqueness argument --- of which the second is blocked in one place that is not a matter of
effort, the regularized incomplete gamma function being absent from Mathlib at the pinned
version, so that dimension $3$ would have to be written out by hand. A second cost the printed
argument hides: a delay datum has no bridge in the development, every consumer so far having had
the monotonicity field, so the five other fields of \ref{eq:sd-profile} would have to be
reproved for it. No node of this article reads any clause of this lemma. Not attempted.)}
\end{lemma}

\begin{proof}
Throughout $m = (d+1)/2$, and $P(s,x) = \Gamma(s)^{-1}\int_0^x v^{s-1}e^{-v}\,dv$ is the
regularized lower incomplete gamma function.

\emph{The datum.} $k^{(c)}_{\mathrm I}$ is nonnegative and measurable, and
\[
  \int_0^1 k^{(c)}_{\mathrm I}(u)\,du = 2, \qquad
  \int_1^\infty k^{(c)}_{\mathrm I}(u)\,\frac{du}{u} = 2 + c\log 2 ,
\]
both finite. So $D^{(c)}$ is a delay datum and $F^{(c)} \in \mathcal{G}$. That
$\mathcal{S} \subseteq \mathcal{G}$ is \ref{def:delay-data} read against
\ref{def:causal-admissible}: a causally admissible datum is a delay datum, its profile being
nonincreasing and hence measurable.

\emph{The profile has no nonincreasing version.} Choose $\delta \in (0,\tfrac12)$ with
$(1-\delta)^{-1/2} - (1+\delta)^{-1/2} < c$, which is possible because the left side tends to $0$
as $\delta$ decreases to $0$. For $u_1 \in (1-\delta,1)$ and $u_2 \in (1,1+\delta)$,
\[
  k^{(c)}_{\mathrm I}(u_1) = u_1^{-1/2} < (1-\delta)^{-1/2} < (1+\delta)^{-1/2} + c
   < u_2^{-1/2} + c = k^{(c)}_{\mathrm I}(u_2) ,
\]
while $u_1 < u_2$. If a nonincreasing $g$ agreed with $k^{(c)}_{\mathrm I}$ almost everywhere then,
both intervals having positive measure, there would be such a pair with
$g(u_i) = k^{(c)}_{\mathrm I}(u_i)$, and $g(u_1) \ge g(u_2)$ would contradict the display. So
$D^{(c)}$ is not a causally admissible datum.

\emph{What the bridge computation reads.} Of the four steps in the proof of
\ref{lem:bridge-exponents}, only the second reads the monotonicity of the delay profile. The
finiteness of $\nu_2(x) = \int_0^\infty g_u(x)k_{\mathrm I}(u)\,du/u$ off the origin uses
$g_u(x)/u \le 16/x^4$ below $u = 1$ and $g_u \le 1$ above it, against the two integrability
conditions; the \Levy{} condition is a Tonelli exchange against
$\int_\RR (1 \wedge x^2)g_u \le 1 \wedge u$ and \ref{lem:profile-integrability}; and the exponent
identity is a second Tonelli exchange. All three read the two integrability conditions and nothing
else, so they hold for a delay datum verbatim. Hence for any delay datum $(b_0, k_{\mathrm I})$
the pair $(b_0/2,\,2x\nu_2(x))$ satisfies every requirement \ref{eq:sd-profile} makes except
possibly the monotonicity of the profile, and its exponent is $F_{\mathrm I}(\omega^2/2)$; so the
datum's exponent is admissible as soon as that profile is nonincreasing. The substitution
$u = r^2/2v$ that produces \ref{eq:bridge-scaling} identifies the profile with $k_1$, the polar
profile \ref{eq:polar-profile} at $d = 1$, the folding factor $2$ being $|S^0|$.

\emph{The polar profiles in closed form.} Split \ref{eq:polar-profile} at the two summands of
$k^{(c)}_{\mathrm I}$. For the power, $k_{\mathrm I}(r^2/2v) = (r^2/2v)^{-1/2}$ turns the integral
into a gamma integral and gives $2^{3/2}r^{-1}\Gamma(m)/\Gamma(d/2)$. For the bump,
$\mathbf{1}_{[1,2]}(r^2/2v) = \mathbf{1}_{[r^2/4,\,r^2/2]}(v)$, so that summand contributes
$2c\,[P(d/2, r^2/2) - P(d/2, r^2/4)]$. Hence
\begin{equation}\label{eq:filtration-polar}
  k^{(c)}_d(r) = \frac{2^{3/2}\,\Gamma(m)}{\Gamma(d/2)}\,\frac1r
   + 2c\Bigl[P\Bigl(\frac d2,\frac{r^2}2\Bigr) - P\Bigl(\frac d2,\frac{r^2}4\Bigr)\Bigr],
   \qquad r > 0 ,
\end{equation}
which is smooth on $(0,\infty)$; the bracket is
$\operatorname{erf}(r/\sqrt2) - \operatorname{erf}(r/2)$ at $d = 1$ and $e^{-r^2/4} - e^{-r^2/2}$
at $d = 2$. Since $\partial_x P(s,x) = x^{s-1}e^{-x}/\Gamma(s)$,
\[
  \frac{\partial}{\partial r}P\Bigl(\frac d2, \frac{r^2}{2a}\Bigr)
   = \frac{2}{r\,\Gamma(d/2)}\Bigl(\frac{r^2}{2a}\Bigr)^{d/2}e^{-r^2/2a} ,
\]
so multiplying the derivative of \ref{eq:filtration-polar} by $r^2\Gamma(d/2)$ and writing
$z = r^2/2$ shows that it is nonpositive at $r$ if and only if
\begin{equation}\label{eq:filtration-criterion}
  c\,z^m\bigl(e^{-z} - 2^{-d/2}e^{-z/2}\bigr) \le \tfrac12\Gamma(m), \qquad z = \tfrac12 r^2 .
\end{equation}

\emph{Two elementary bounds.} First, $\log(1/p) \le (1-p)/\sqrt p$ for $0 < p \le 1$: with
$x = 1/p \ge 1$ this reads $\log x \le (x-1)/\sqrt x$, an equality at $x = 1$ where the right side
has the larger derivative from then on, since $1/x \le (x+1)/(2x^{3/2})$ is
$2\sqrt x \le x + 1$. Second, for $\lambda \in (0,1)$ and $m > 0$ the function
$(1-p)^m(p-\lambda)$ on $[\lambda,1]$ is maximal at $p = (m\lambda+1)/(m+1)$, with value
$m^m(1-\lambda)^{m+1}/(m+1)^{m+1}$.

\emph{Where the polar profile is nonincreasing.} Put $p = e^{-z/2}$ and $\lambda = 2^{-d/2}$, so
that $e^{-z} - \lambda e^{-z/2} = p(p-\lambda)$, which is positive only for $p > \lambda$, and
there $z = 2\log(1/p) \le 2(1-p)/\sqrt p$. So for $m \le 2$, using $p^{1-m/2} \le 1$ and then the
second bound,
\[
  z^m\bigl(e^{-z} - \lambda e^{-z/2}\bigr) \le 2^m(1-p)^m p^{1-m/2}(p-\lambda)
   \le 2^m(1-p)^m(p-\lambda) \le \frac{(2m)^m(1-\lambda)^{m+1}}{(m+1)^{m+1}} .
\]
At $d = 1$, where $m = 1$, $\lambda = 2^{-1/2}$ and $\Gamma(m) = 1$, the right side is
$(1-2^{-1/2})^2/2 = (\tfrac32-\sqrt2)/2$, so \ref{eq:filtration-criterion} holds at every $z$ as
soon as $c(\tfrac32-\sqrt2) \le 1$, that is as soon as $c \le 6 + 4\sqrt2$, the two numbers being
reciprocal: $(\tfrac32-\sqrt2)(6+4\sqrt2) = 9 - 8 = 1$. At $d = 2$, where $m = \tfrac32$,
$\lambda = \tfrac12$ and $\Gamma(m) = \sqrt\pi/2$, the right side is
$3^{3/2}5^{-5/2} = (3/5)^{3/2}/5$, so at $c = 4$ the criterion holds at every $z$ as soon as
$4\cdot 3^{3/2}5^{-5/2} \le \sqrt\pi/4$, that is $(16/5)(3/5)^{3/2} \le \sqrt\pi$, that is, after
squaring, $6912/3125 \le \pi$, which holds, the left side being less than $2.22$. Since $4$ and
$8$ are both below $6+4\sqrt2$, this gives that $k^{(c)}_1$ is nonincreasing for
$0 < c \le 6+4\sqrt2$ and that $k^{(4)}_2$ is nonincreasing.

\emph{Where it is not.} If \ref{eq:filtration-criterion} fails at one $z_0 > 0$ then $k^{(c)}_d$
has a positive derivative at $r_0 = \sqrt{2z_0}$, so it is strictly increasing near $r_0$; being
continuous, it agrees with no nonincreasing function almost everywhere either. At $d = 2$ and
$c = 8$ take $z_0 = \tfrac12$. Then $e^{-z_0} > 3/5$ and $e^{-z_0/2} < 4/5$ --- the first is
$e < 25/9$, the second $e > (5/4)^4 = 625/256$ --- so
$e^{-z_0} - \tfrac12 e^{-z_0/2} > 3/5 - 2/5 = 1/5$; and $z_0^{3/2} = 2^{-3/2} > 7/20$, since
$400 > 392$. The left side of \ref{eq:filtration-criterion} therefore exceeds
$8 \cdot \tfrac7{20} \cdot \tfrac15 = \tfrac{14}{25}$, while its right side $\sqrt\pi/4$ is below
$0.45$ because $\pi < 3.24$. At $d = 3$ and $c = 4$ take $z_0 = 1$, where
$\Gamma(m) = \Gamma(2) = 1$. Then $e^{-1} > 9/25$, again because $e < 25/9$, while
$2^{-3/2} < 9/25$ because $625 < 648$ and $e^{-1/2} < 61/100$ because $e > (100/61)^2$. So
\[
  e^{-1} - 2^{-3/2}e^{-1/2} > \frac9{25} - \frac9{25}\cdot\frac{61}{100} = \frac{351}{2500} ,
\]
and the left side of \ref{eq:filtration-criterion} exceeds $4\cdot 351/2500 > \tfrac12$. So
$k^{(8)}_2$ and $k^{(4)}_3$ are not nonincreasing, which completes clause (2).

\emph{No causal ancestor.} Suppose $F^{(c)} \in \mathcal{S}$. By \ref{prop:bridge-strictness}(1)
the function $\sigma \mapsto F^{(c)}(\sqrt{2\sigma}) = F^{(c)}_{\mathrm I}(\sigma)$ is causally
admissible, so there are $b_0' \ge 0$ and a nonincreasing $k'$ satisfying the two integrability
conditions whose exponent is $F^{(c)}_{\mathrm I}$ on $[0,\infty)$. The differentiation step in the
proof of \ref{prop:bridge-strictness}(2) dominates the $\sigma$-derivative of the integrand by
$e^{-\sigma u/2}k_{\mathrm I}(u)$ and reads the two integrability conditions and nothing else, so
it applies to both data and gives
\[
  b_0' + \int_0^\infty e^{-\sigma u}k'(u)\,du
   = \int_0^\infty e^{-\sigma u}k^{(c)}_{\mathrm I}(u)\,du , \qquad \sigma > 0 .
\]
Both integrals are finite and tend to $0$ as $\sigma$ grows: for $\sigma \ge 1$ the integrand is
at most $k_{\mathrm I}$ on $(0,1)$ and at most $k_{\mathrm I}(u)u^{-1}$ on $(1,\infty)$, since
$e^{-u} \le u^{-1}$ there, and it tends to $0$ at every $u > 0$, so dominated convergence applies.
Hence $b_0' = 0$, and the measures $k'(u)\,du$ and $k^{(c)}_{\mathrm I}(u)\,du$ on $(0,\infty)$
have equal finite Laplace transforms on $(0,\infty)$. By
\ref{prop:laplace-uniqueness-locally-finite} they are equal, so $k'$ is a nonincreasing version of
$k^{(c)}_{\mathrm I}$, which the second step excludes. So $F^{(c)} \notin \mathcal{S}$ for every
$c > 0$; for $0 < c \le 6+4\sqrt2$ the fifth step makes $F^{(c)}$ admissible with Gaussian
coefficient $0$ and profile $k^{(c)}_1$, and clause (1) follows.

\emph{The polar profile is the polar profile.} Clause (3) reads clause (2) through the polar
criterion, and the reading needs two steps. The first is that $k^{(c)}_d$ is the polar profile of
a law on $\RR^d$ with radial exponent $\xi \mapsto F^{(c)}(|\xi|)$, which is the computation of
\ref{prop:dimension-filtration}(2) carried out at a datum that need not be monotone. That
computation builds $\nu_d(dx) = \bigl(\int_0^\infty g^{(d)}_u(x)\,k_{\mathrm I}(u)\,du/u\bigr)dx$
and establishes, in this order: that $\nu_d$ is finite off the origin and satisfies the \Levy{}
condition, by a Tonelli exchange against
$\int_{\RR^d}(1 \wedge |x|^2)g^{(d)}_u \le 1 \wedge du$; that
$F_{\mathrm I}(|\xi|^2/2) = \tfrac12\xi\cdot(b_0I)\xi + \int(1 - \cos\xi\cdot x)\,\nu_d(dx)$, by a
second Tonelli exchange; that this is therefore the exponent of an infinitely divisible law,
isotropic because its transform is, whose \Levy{} measure is $\nu_d$; and that the polar profile
of $\nu_d$, namely $|S^{d-1}|r^df_d(r)$ with $f_d$ the radial density, is \ref{eq:radial-profile},
by the substitution $u = r^2/2v$. Each of those steps reads the two integrability conditions of
the datum and nothing else --- the first two are integral bounds, the third is the converse clause
of the representation in dimension $d$ with its uniqueness clause, and the fourth is a change of
variables. Monotonicity of $k_{\mathrm I}$ enters only in the last line of that proof, where the
polar profile is concluded to be nonincreasing. So the computation applies verbatim to the delay
datum $D^{(c)}$, and gives an isotropic infinitely divisible law $\mu^{(c)}_d$ on $\RR^d$ with
radial exponent $F^{(c)}(|\xi|)$ and \Levy{} measure of polar profile $k^{(c)}_d$,
\ref{eq:polar-profile} being \ref{eq:radial-profile}.

\emph{A rotation-invariant \Levy{} measure of a self-decomposable law has a nonincreasing radial
profile.} This is the second step, and it is the polar criterion averaged over the sphere. Let
$\mu$ be self-decomposable on $\RR^d$ with \Levy{} measure $\nu$. The criterion gives a finite
measure $\lambda$ on the unit sphere and, for $\lambda$-almost every direction $\xi$, a
nonincreasing $k_\xi : (0,\infty) \to [0,\infty)$ with
$\nu(B) = \int_{S^{d-1}}\lambda(d\xi)\int_0^\infty \mathbf 1_B(r\xi)\,k_\xi(r)\,dr/r$. Take $B$ of
the form $\{x : |x| \in A\}$ for a Borel $A \subseteq (0,\infty)$ and put
$K(r) := \int_{S^{d-1}}k_\xi(r)\,\lambda(d\xi)$, a function with values in $[0,\infty]$ that is
nonincreasing, each integrand being so. Then $\nu\{|x| \in A\} = \int_A K(r)\,dr/r$ for every such
$A$. If moreover $\nu$ is rotation invariant with radial density $f$, the same quantity is
$|S^{d-1}|\int_A f(r)\,r^{d-1}dr = \int_A |S^{d-1}|r^df(r)\,dr/r$. Two measures on $(0,\infty)$
that agree on every Borel set are equal, so $|S^{d-1}|r^df(r) = K(r)$ for almost every $r$: the
polar profile agrees almost everywhere with a nonincreasing function.

\emph{Clause (3).} Both $F^{(4)}$ and $F^{(8)}$ are admissible, $4$ and $8$ being below
$6 + 4\sqrt2$, so both lie in $\mathcal{A}_1$, which is the whole admissible cone
(\ref{prop:dimension-filtration}(1)). For $F^{(4)} \in \mathcal{A}_2$: the first step gives an
infinitely divisible law on $\RR^2$ with radial exponent $F^{(4)}(|\xi|)$ whose \Levy{} measure
has polar profile $k^{(4)}_2$, which clause (2) says is nonincreasing, so the law is
self-decomposable by the polar criterion. For the two exclusions, suppose
$\xi \mapsto F^{(c)}(|\xi|)$ were the exponent of a self-decomposable law $\mu$ on $\RR^d$, with
$(c,d) = (8,2)$ or $(4,3)$. An exponent determines the law, so $\mu$ is the $\mu^{(c)}_d$ of the
first step, whose \Levy{} measure is rotation invariant with radial density $f^{(c)}_d$ and polar
profile $k^{(c)}_d$. By the second step $k^{(c)}_d$ agrees almost everywhere with a nonincreasing
function; but $k^{(c)}_d$ is continuous and, by the paragraph \emph{where it is not} above,
strictly increasing on an interval, so it agrees with no nonincreasing function almost everywhere
--- given $r < s$, points $u_n \downarrow r$ and $v_n \uparrow s$ of the conull agreement set give
$k^{(c)}_d(u_n) \ge k^{(c)}_d(v_n)$, and continuity would force $k^{(c)}_d(r) \ge k^{(c)}_d(s)$ at
every such pair. That contradiction gives $F^{(8)} \notin \mathcal{A}_2$ and
$F^{(4)} \notin \mathcal{A}_3$. The last sentence of the clause is
$\mathcal{A}_\infty \subseteq \mathcal{A}_3$, which is the nesting of
\ref{prop:dimension-filtration}(1).
\end{proof}

% CHANGED (2026-09-15, module B step 2, R161): the closing CHECK is rewritten. Its question is
% still open in the literature, and the round that read the held sources is recorded there: no
% source settles it, Bondesson settles the analogue one class up, and a construction that answers
% it in the negative is written out at rem:dimension-filtration and is not yet a node. The remark's
% own text is untouched. (2026-09-17: the construction is lem:filtration-strictness since R163, and
% the filtration itself is prop:dimension-filtration since R166.)
\begin{remark}[what crosses the bridge, and what does not]\label{rem:bridge-crossings}
Three of the causal article's later results have images and one has none. The \emph{exact
cascade implementation} crosses: the rational Laplace transform of the causal Gamma increments
becomes the rational spectrum of the \Matern{} increments, and the causal pole--zero sections
become forward--backward first-order sections in $x$ (\ref{prop:matern-cascade}). The \emph{cone
geometry} crosses as a map of cones but not as a map of extreme rays, by
\ref{prop:bridge-strictness}: the image of a causal Dickman ray has the smooth positive profile
of \ref{prop:thorin-subclass}(4), not a step. The \emph{moment and tail dictionary} crosses with
a change of exponent, since $\mathbb{E}|\BM_T|^{2n} = (2n-1)!!\;\mathbb{E}T^n$, so the spatial
family has $2n$ finite moments exactly when the causal one has $n$. What does not cross as an
\emph{object} is the memory line: the causal signaling form and its embodied jet live on the
inversion operator of the memory half-line, which exists because of the boundary and embodiment
structure of time, and the spatial theory has no such second operator. % CHANGED 2026-09-19 (R195, the author's decision at the second review round): the \ref into
% chapter 12 is replaced by prose. Module C is not published, and a reference from a module-B
% chapter into a module-C statement would break at a repository split; the mathematics of the
% sentence is unchanged and what it announces is named as further work.
% TODO(module C): cite thm:joint-locality here once module C is released.
What does cross is the
memory line's \emph{equation}: the causal locality theorem, that the memory-line evolution is a
local second-order operator exactly for the Bessel family, has as its image here a statement that
the Student-t scale space solves a local second-order equation in $(t,x)$ jointly, with the same
Bessel operator now acting along the scale axis. That statement is left for further work and is
made nowhere in this module; what belongs to this chapter is the substitution
$\sigma \mapsto \omega^2/2$, which turns $\partial_t$ into $-\partial_x^2$, a parabolic equation
into an elliptic one. So the causal local corner maps to the spatial jointly local corner; the
asymmetry that remains is that the spatial theory has no signaling problem and no state to read
a jet off, and that, not any failure of the bridge, is why the two articles' second halves
differ.
% CHECK (draft, checks for \S9, item 7; the literature round is 2026-09-15, R161): whether the
% subordinated slice coincides with the symmetric self-decomposable laws of type G (variance
% mixtures of normals with infinitely divisible mixing) or is strictly smaller (mixing law
% self-decomposable) is OPEN IN THE HELD LITERATURE, and the sources are now read, so the
% literature paragraph can be written either way. What they give, with pages, is in
% rem:dimension-filtration: Halgreen 1979 pp. 14-16 proves the implication in the direction the
% bridge uses and leaves a DIFFERENT question open (the mean-variance case); Barndorff-Nielsen and
% Halgreen 1977 pp. 309-311 give infinite divisibility of the mixture in every dimension;
% Bondesson 1992 Thm. 7.3.1 pp. 115-116 settles the analogue one class up, at the Thorin level, as
% an EQUIVALENCE. None states the converse at the level of class L.
% What the same round found and did not state as a node: an explicit admissible exponent of type G
% whose mixing subordinator is not self-decomposable, which answers the question in the negative.
% It is the closing paragraph of rem:dimension-filtration and SKELETON.md section 23, and it is a
% candidate lemma for the author. Until the author decides, the CHECK stands and nothing in this
% chapter asserts an answer; no statement here reads it either way.
\end{remark}
